% !TeX program = xelatex
% CS224 课程项目最终报告（中文，IEEE 会议格式）。
% 结果来自 report/generated_tables.txt、report/generated_costs.json、analysis/aggregation.json、
% qualitative/cases.json 与 repair_eval/analysis/repair_blind_v1/final/；不编造数值。
\RequirePackage[T1]{fontenc}
\documentclass[conference]{IEEEtran}

% scheme=plain 保留 IEEEtran 的标题、章节和页面设置。
\usepackage[UTF8,scheme=plain,fontset=fandol]{ctex}
\setmainfont{Nimbus Roman}
\setsansfont{Nimbus Sans}
\setmonofont{Nimbus Mono PS}
\setCJKmainfont{FandolSong-Regular.otf}[
    BoldFont=FandolSong-Bold.otf,
    ItalicFont=FandolKai-Regular.otf,
    BoldItalicFont=FandolKai-Regular.otf,
    SmallCapsFont=FandolSong-Regular.otf]
\usepackage{cite}
\usepackage{amsmath,amssymb}
\usepackage{booktabs}
\usepackage{graphicx}
\usepackage{multirow}
\usepackage{array}
\usepackage{xcolor}
\usepackage{listings}
\usepackage{balance}
\usepackage{xurl}
\usepackage[hidelinks]{hyperref}
\urlstyle{tt}
\hypersetup{pdfauthor={Yu Zhang},pdftitle={C/C++ 漏洞检测、配对判别与修复的 LLM 评估：统一分层基准、双模型比较与编译闭环}}
\renewcommand{\abstractname}{摘要}
\renewcommand{\IEEEkeywordsname}{关键词}
\renewcommand{\refname}{参考文献}
\renewcommand{\thetable}{\arabic{table}}

\lstset{
  basicstyle=\ttfamily\scriptsize,
  breaklines=true,
  columns=fullflexible,
  keepspaces=true,
  showstringspaces=false,
  frame=single,
  rulecolor=\color{black},
  xleftmargin=0.5em,
  xrightmargin=0.5em,
  aboveskip=0.5em,
  belowskip=0.5em
}


\newcommand{\Dind}{\mathcal{D}_{\mathrm{ind}}}
\newcommand{\Dpair}{\mathcal{D}_{\mathrm{pair}}}

\title{C/C++ 漏洞检测、配对判别与修复的 LLM 评估：\\统一分层基准、双模型比较与编译闭环}
\author{
\IEEEauthorblockN{Yu Zhang}
\IEEEauthorblockA{CS224 课程项目\\
ShanghaiTech University\\
Shanghai, China\\
\href{mailto:Zhangyu12026@shanghaitech.edu.cn}{Zhangyu12026@shanghaitech.edu.cn}}
% TODO: 按课程分组补充全部作者、单位与邮箱。
}

% Auto-generated by report/scripts/build_extended_tables.py. Do not edit.
\newcommand{\NumTADSDirMF}{0.383}
\newcommand{\NumTADSFewMF}{0.399}
\newcommand{\NumTADSConMF}{0.381}
\newcommand{\NumTAMiDirMF}{0.346}
\newcommand{\NumTAMiFewMF}{0.346}
\newcommand{\NumTAMiConMF}{0.365}
\newcommand{\NumTADSDirCILo}{0.305}
\newcommand{\NumTADSDirCIHi}{0.552}
\newcommand{\NumTADSFewCILo}{0.305}
\newcommand{\NumTADSFewCIHi}{0.593}
\newcommand{\NumTADSConCILo}{0.286}
\newcommand{\NumTADSConCIHi}{0.560}
\newcommand{\NumTAMiDirCILo}{0.264}
\newcommand{\NumTAMiDirCIHi}{0.528}
\newcommand{\NumTAMiFewCILo}{0.262}
\newcommand{\NumTAMiFewCIHi}{0.512}
\newcommand{\NumTAMiConCILo}{0.278}
\newcommand{\NumTAMiConCIHi}{0.571}
\newcommand{\NumTADSPairDir}{0.117}
\newcommand{\NumTADSPairFew}{0.089}
\newcommand{\NumTADSPairCon}{0.122}
\newcommand{\NumTAMiPairDir}{0.091}
\newcommand{\NumTAMiPairFew}{0.094}
\newcommand{\NumTAMiPairCon}{0.114}
\newcommand{\NumTAModelFewDelta}{0.053}
\newcommand{\NumTAModelFewCILo}{0.010}
\newcommand{\NumTAModelFewCIHi}{0.129}
\newcommand{\NumTAModelFewRawP}{0.0784}
\newcommand{\NumTAModelFewHolmP}{0.2351}
\newcommand{\NumTADSCoverage}{1.000}
\newcommand{\NumTAMiCoverage}{0.945}
\newcommand{\NumTADSFullUSD}{2.439219}
\newcommand{\NumTAMiFullUSD}{0.623154}
\newcommand{\NumTADSUSDPerHundred}{0.098197}
\newcommand{\NumTAMiUSDPerHundred}{0.025087}
\newcommand{\NumPJDSDA}{76.67}
\newcommand{\NumPJMiDA}{86.67}
\newcommand{\NumPJDSA}{92.22}
\newcommand{\NumPJDSB}{61.11}
\newcommand{\NumPJMiA}{88.89}
\newcommand{\NumPJMiB}{84.44}
\newcommand{\NumPJDSFishP}{\ensuremath{1.00\times10^{-6}}}
\newcommand{\NumPJMiFishP}{0.5114}
\newcommand{\NumPJHeur}{74.44}
\newcommand{\NumPJDSLenP}{0.6587}
\newcommand{\NumPJMiLenP}{\ensuremath{3.13\times10^{-4}}}
\newcommand{\NumPJDSUSD}{0.153593}
\newcommand{\NumPJMiUSD}{0.071803}
\newcommand{\NumEvDSRecBase}{29.1}
\newcommand{\NumEvDSRec}{38.6}
\newcommand{\NumEvDSDeltaRec}{9.4}
\newcommand{\NumEvDSFPRBase}{25.0}
\newcommand{\NumEvDSFPR}{34.4}
\newcommand{\NumEvDSDeltaFPR}{9.4}
\newcommand{\NumEvDSStrictBase}{38.8}
\newcommand{\NumEvDSStrict}{43.8}
\newcommand{\NumEvDSPPRBase}{28.3}
\newcommand{\NumEvDSPPR}{37.7}
\newcommand{\NumEvDSUSDBase}{0.066988}
\newcommand{\NumEvDSUSD}{0.085622}
\newcommand{\NumEvMiRecBase}{24.4}
\newcommand{\NumEvMiRec}{35.3}
\newcommand{\NumEvMiDeltaRec}{10.9}
\newcommand{\NumEvMiFPRBase}{20.7}
\newcommand{\NumEvMiFPR}{34.5}
\newcommand{\NumEvMiDeltaFPR}{13.8}
\newcommand{\NumEvMiStrictBase}{32.5}
\newcommand{\NumEvMiStrict}{41.9}
\newcommand{\NumEvMiPPRBase}{23.6}
\newcommand{\NumEvMiPPR}{35.1}
\newcommand{\NumEvMiUSDBase}{0.024397}
\newcommand{\NumEvMiUSD}{0.035933}
\newcommand{\NumEvDSHolmRec}{0.0055}
\newcommand{\NumEvDSHolmFPR}{0.2500}
\newcommand{\NumEvMiHolmRec}{0.0010}
\newcommand{\NumEvMiHolmFPR}{0.2500}
\newcommand{\NumEvDSExact}{0.488}
\newcommand{\NumEvDSChance}{0.345}
\newcommand{\NumEvMiExact}{0.455}
\newcommand{\NumEvMiChance}{0.316}
\newcommand{\NumJbClangMF}{0.888}
\newcommand{\NumJbClangCILo}{0.818}
\newcommand{\NumJbClangCIHi}{0.948}
\newcommand{\NumJbClangMCC}{0.799}
\newcommand{\NumJbClangCov}{0.971}
\newcommand{\NumJbClangStrict}{0.864}
\newcommand{\NumJbCppMF}{0.588}
\newcommand{\NumJbCppCILo}{0.497}
\newcommand{\NumJbCppCIHi}{0.671}
\newcommand{\NumJbCppMCC}{0.331}
\newcommand{\NumJbCppCov}{1.000}
\newcommand{\NumJbCppStrict}{0.629}
\newcommand{\NumJbDynMF}{0.914}
\newcommand{\NumJbDynCILo}{0.850}
\newcommand{\NumJbDynCIHi}{0.966}
\newcommand{\NumJbDynMCC}{0.841}
\newcommand{\NumJbDynCov}{1.000}
\newcommand{\NumJbDynStrict}{0.914}
\newcommand{\NumJbDynConsN}{58}
\newcommand{\NumJbDynConsCov}{0.414}
\newcommand{\NumJbDynConsAbstain}{82}
\newcommand{\NumJlAcc}{95.00}
\newcommand{\NumJlF}{0.951}
\newcommand{\NumJlMCC}{0.901}
\newcommand{\NumJlRepAcc}{95.00}
\newcommand{\NumJlRepF}{0.951}
\newcommand{\NumJlCost}{0.003180}
\newcommand{\NumJlPreAcc}{95.71}
\newcommand{\NumJlPreF}{0.958}
\newcommand{\NumJlCAcc}{93.10}
\newcommand{\NumJlCppAcc}{96.34}
\newcommand{\NumHyDSDir}{96.43}
\newcommand{\NumHyDSDyn}{97.14}
\newcommand{\NumHyDSComb}{97.14}
\newcommand{\NumHyDSDelta}{0.71}
\newcommand{\NumHyMiDir}{97.14}
\newcommand{\NumHyMiStatic}{96.43}
\newcommand{\NumHyMiComb}{97.14}
\newcommand{\NumHyMiDelta}{0.00}
\newcommand{\NumHyDSCombP}{1.0000}
\newcommand{\NumHyMiCombP}{1.0000}
\newcommand{\NumHyTotalUSD}{0.082335}
\newcommand{\NumHyTotalCalls}{1400}
\newcommand{\NumHyChangedUnits}{6}
\newcommand{\NumHyFlipRightWrong}{16}
\newcommand{\NumHyFlipWrongRight}{2}
\newcommand{\NumHyVoneDSComb}{100.00}
\newcommand{\NumHyVoneDyn}{99.29}
\newcommand{\NumHyPooledVone}{97.36}
\newcommand{\NumHyPooledVtwo}{96.36}
\newcommand{\NumCostHeadlineConf}{0.558073}
\newcommand{\NumCostHeadlineNorm}{0.365704}
\newcommand{\NumCostAllConf}{3.886970}
\newcommand{\NumCostAllNorm}{2.146513}
\newcommand{\NumCostAllRuns}{71}
\newcommand{\NumCostAllCalls}{9,532}
\newcommand{\NumCostRepairHeadConf}{0.032831}
\newcommand{\NumCostRepairHeadNorm}{0.019336}
\newcommand{\NumCostRepairAllConf}{0.082550}
\newcommand{\NumCostRepairAllNorm}{0.052549}
\newcommand{\NumCostRepairOld}{0.077516}
\newcommand{\NumCostRepairDelta}{0.005034}
\newcommand{\NumRepObs}{58}
\newcommand{\NumRepUnits}{70}
\newcommand{\NumRepDSStrict}{58}
\newcommand{\NumRepDSFirst}{58}
\newcommand{\NumRepDSCalls}{70}
\newcommand{\NumRepDSCost}{0.026990}
\newcommand{\NumRepMiStrict}{58}
\newcommand{\NumRepMiFirst}{57}
\newcommand{\NumRepMiCalls}{72}
\newcommand{\NumRepMiCost}{0.005841}
\newcommand{\NumRepTotalCost}{0.032831}


\begin{document}
\maketitle

\begin{abstract}
大语言模型（large language model，LLM）已被广泛用于源代码漏洞检测与修复，但不同提示策略、不同模型以及不同评估层级的结果，常因标签、分母、成本和失败口径不同而无法并列阅读。本文在 SecVulEval 冻结测试集（160 条检测函数、180 对漏洞/修复版本）上比较零样本直接判定、固定少样本提示和白名单上下文增强三种策略，并把实验扩展为四个真实部分：DeepSeek 与 MiMo 的双模型复制、180 对强制二选一联合提示、160 条结构化证据提示，以及 Juliet CWE-401 的静态/动态/LLM/证据混合与修复编译闭环。Tier A 检测 macro-F1 为 DeepSeek 0.383/0.399/0.381、MiMo 0.346/0.346/0.365；6 个同策略模型比较经 Holm 校正后均不显著；MiMo few-shot 的差值区间 $[0.010,0.129]$ 虽排除 0，但 McNemar Holm $p=0.235$，只能视为提示性。分侧配对 DA 为 8.9\%--12.2\%（all-pairs 7.8\%--12.2\%），远低于 56.1\% 的机会水平。强制二选一把 DA 提到 76.7\%/86.7\%，但“选较短版本”的长度启发式单独达到 74.4\%，DeepSeek 与启发式不可区分（$p=0.659$），且存在 31.1 个百分点的 A 位置偏差；只有 MiMo 显著高于长度启发式（$p=0.0003$）。证据提示把召回提高 9.4/10.9 个百分点（Holm $p=0.0055/0.0010$），但 FPR 同时上升 9.4/13.8 个百分点，平衡准确率与 MCC 基本持平，收益主要是操作点移动。Juliet test 140 上 Clang SA macro-F1 0.888、Cppcheck 0.588、ASan/Valgrind 0.914；动态 clean 只表示未观察到泄漏，不是 TN。MiMo 在 Juliet 上（泄漏修复后）为 95.0\%（F1 0.951，Xiaomi 官方端点 fallback，非 OpenCode Go）；证据混合 v2 的最好单元为 97.1\%/97.1\%，相对 direct 的 $+1/0$ 个单元在 McNemar 下不显著，安慰剂无聚合收益。修复 6 个 good 单元的标签字面量后，DeepSeek combined 从 100\% 回落到 97.14\%；v1 的 100\% 只作审计案例，不作为主结论。修复闭环中两个模型在 58 个 LSan 可观测单元上均 58/58 严格通过，12 个不可观测单元不计入。结论强调：统一契约、分层报告和显式失败/混杂审计比单一 leaderboard 更重要；所有结果均为单次运行，不声称稳定优劣或通用修复正确率。
\end{abstract}

\begin{IEEEkeywords}
漏洞检测，修复生成，大语言模型，提示策略，分层基准，静态分析，动态分析，可复现性
\end{IEEEkeywords}

\section{引言}
\label{sec:intro}
源代码中的安全缺陷是软件供应链风险的重要来源。传统上，漏洞发现依赖人工审计、模糊测试和静态分析规则；这些手段各有覆盖范围，但面对大型代码库时成本较高。基于学习的漏洞检测试图从已标注的函数或代码片段中学习模式，从而为开发者提供可扩展的候选告警。早期工作以代码切片、图神经网络或预训练语言模型为核心，在函数级二分类任务上取得了可观进展~\cite{li2018vuldeepecker,zhou2019devign,feng2020codebert,fu2022linevul}。

近年来，通用大语言模型展现出跨语言、跨任务的代码理解能力，使“不训练、只提示”的漏洞判定与修复生成成为可能~\cite{chen2021codex,pearce2023zeroshot}。然而，评估这类方法并不只是报告一个准确率。首先，漏洞样本与无漏洞样本的比例会显著影响精确率、召回率和 F1 的解释；其次，代码模型可能在预训练阶段见过公开漏洞及其修复，导致评估高估泛化能力；再次，单函数分类无法回答两个更接近实践的问题：当同一个函数被修复后，检测器能否同时判对两个版本？模型能否生成一个完整、与输入不同、且与参考修复接近的补丁？前者衡量“修复差异辨识”，后者衡量“修复生成”，二者不能互相替代~\cite{ding2025primevul,pearce2023zeroshot}。

SecVulEval 为上述问题提供了合适的数据基础。该基准面向 LLM 的上下文感知漏洞检测，包含 5{,}867 个 CVE、10{,}998 条有漏洞函数和 14{,}442 条无漏洞函数，并提供语句级变更标注与函数上下文~\cite{ahmed2026secvuleval,secvulevaldata}。本项目从本地 SecVulEval 快照中抽取两个互补源子集：一个用于函数级检测的独立集合，另一个由同一函数的漏洞版本与修复版本构成的配对集合。两个源子集在原始索引和“项目—提交”分组上相互隔离，合计 600 条不同的函数记录。随后按项目级隔离冻结 \texttt{secvul-split-v1}：检测集 40 条开发（32 正/8 负）+ 160 条测试（128 正/32 负），配对集 20 对开发 + 180 对测试。模型评估只使用测试集，开发样本只用于固定少样本示例和 pilot。源子集不是训练集，而是两种评估视图；冻结划分决定哪些记录进入主结果。

本文关注三种可在相同基础模型上复现的 LLM 提示策略：零样本直接判定、固定少样本提示，以及结合数据集白名单上下文的增强提示。三者共享任务描述、输出格式、解析器和计费口径，差异仅在提示中提供的信息。这样可以把“策略带来的差异”与“模型本身的能力差异”分开。本文的研究问题如下。

\begin{enumerate}
    \item \textbf{RQ1（函数级检测）}：在 160 条检测测试函数上，三种策略的精确率、召回率、F1、macro-F1、特异度、平衡准确率、MCC 和无效输出率分别是多少？
    \item \textbf{RQ2（配对联合判别）}：在 180 对测试版本上，三种策略能否同时判对同一函数的两个版本？联合判别准确率和四类联合结果（正确区分、双报漏洞、双报安全、漏报加误报）分别是什么？
    \item \textbf{RQ3（修复生成）}：在 308 条测试漏洞版本上，三种策略生成的修复在模式有效、函数完整、与输入不同、完整且不同、与参考修复文本一致等结构指标上表现如何？
    \item \textbf{RQ4（成本）}：三种策略的 token 消耗、费用与延迟差异有多大？每 100 个评估单元和每正确预测的成本是多少？
    \item \textbf{RQ5（误差类型）}：三种策略在假阳性、假阴性、配对失败和修复失败上呈现哪些可归类的模式？
    \item \textbf{RQ6（跨模型复制）}：把同一套提示和冻结划分复制到 MiMo-V2.6-Flash 后，检测、配对和成本结论是否与 DeepSeek 一致？同策略、同任务的模型差异在 Holm 校正后是否显著？
    \item \textbf{RQ7（联合提示）}：把两次分侧判定改成每对一次强制二选一，配对 DA 能提高多少？提高中有多少可由位置偏差和“修复版本通常更长”的长度启发式解释？
    \item \textbf{RQ8（证据与分层）}：结构化证据行/反事实提示能否提高检测召回而不增加假阳性？在 Juliet CWE-401 上，静态、动态、LLM 和证据混合各自达到什么水平，哪些差异在分层分母下仍可解释？
    \item \textbf{RQ9（修复闭环）}：在真实编译、LeakSanitizer 和行为差分 oracle 下，两个模型能修复多少个 test bad 单元？不可观测单元应如何从分母中排除？
\end{enumerate}

需要先说明研究范围。项目最初设想聚焦 C++ 内存泄漏或 Juliet CWE-401 专项，但本地源子集是跨 CWE 的真实项目抽样，不是按 CWE-401 构造的基准：360 条源漏洞函数中只有 4 条带 CWE-401 标签（独立集 2/160，配对集 2/200；冻结测试集 308 条中 3 条，唯一一条 C++ 记录 17743 落在开发集），其中仅 1 条为 C++ 扩展名且同时带 CWE-415，严格“C++ 扩展名且单一 CWE-401 标签”为 0 条；两个 JSON 中也没有 Juliet 记录。若按 CWE-401 专项报告指标，任何数值都会由 1--2 条样本决定。因此本文把主问题定位为\textbf{全 CWE} 的函数级检测与漏洞/修复联合判别，CWE 只作描述性分层，CWE-401 只作案例系列与长尾覆盖限制，不作专项定量结论，也不把结果外推到 Juliet CWE-401。扩展部分另行构建 Tier B Juliet CWE-401 子集，只在该层内报告静态、动态、LLM 与证据混合结果；它与 Tier A 的 SecVulEval 数据、标签、分母和成本口径都不同，二者不池化、不横比。Tier A 的“全 CWE 真实项目函数”与 Tier B 的“单 CWE 合成用例”是互补的压力测试，而不是同一排行榜的两行。

本文的贡献是：
\begin{itemize}
    \item 在 SecVulEval 上给出一个可复现的双任务评估协议，覆盖独立函数检测、配对联合判别和修复生成，并明确字段屏蔽、输出模式与评分规则；
    \item 在固定基础模型下并列比较三种提示策略，统一解析、重试、计费和统计不确定性口径；
    \item 给出检测指标、配对联合指标、修复结构代理、成本与误差类型的联合分析框架，并公开提示模板、输出模式与逐样本日志格式；
    \item 给出 Tier A/Tier B 的统一分层契约：同一套字段、失败分类、coverage/strict 和成本向量，但两个层级各自设主终点与外推边界，拒绝跨层池化；
    \item 把 DeepSeek 的三策略实验复制到 MiMo-V2.6-Flash，共 12 个 Tier A run、每模型 2{,}484 次请求，并用同策略 McNemar + Holm 与 cluster bootstrap 报告模型差异；
    \item 设计并真实运行 180 对强制二选一联合提示，配合 HMAC 平衡的 A/B 顺序、位置偏差检验和事后的长度启发式审计；
    \item 设计并真实运行 160 条结构化证据提示（证据行、机制、反事实修复、自检），并把证据行与 \texttt{changed\_lines} 及随机引用基线比较；
    \item 在 Juliet CWE-401 上完成静态、动态、LLM 和证据混合四类真实运行，其中混合实验包含 2 个模型 $\times$ 5 种方法 $\times$ 140 个单元、共 1{,}400 次请求和错配证据安慰剂；发现 6 个 good 单元含标签字面量后，用修复后的 v2 输入包重跑全部 1{,}400 次请求，v1 仅保留为审计对照；
    \item 完成 70 个 Juliet test bad 单元的修复闭环：模型输出经过真实编译、LSan 和行为差分，严格通过只统计 baseline 确实可观测泄漏的 58 个单元。
\end{itemize}

本文报告的是冻结测试集上每个策略的一次完整运行；不以估计值、文献值或 mock 输出代替实测结果。由于每个模型/策略单元只有一次运行，本文不把策略间的小差异解释为稳定优势，配对与修复结论均按探索性证据处理。扩展部分的 MiMo 运行、联合提示、证据提示、Juliet 基准和修复闭环各自有独立的运行次数、分母和限制，均在对应章节单独说明；读者不应把不同章节的点估计直接相减。

\section{相关工作}
\label{sec:related}

\subsection{基于学习的漏洞检测}
VulDeePecker 将函数切片为语义相关的代码 gadget，并用双向循环网络判断切片是否包含漏洞~\cite{li2018vuldeepecker}。Devign 把函数表示为融合抽象语法、控制流和数据流的图，再用图神经网络学习函数级表示~\cite{zhou2019devign}。这类方法强调程序结构，但依赖大规模标注数据，且跨项目泛化常受数据偏差影响。

ReVeal 重新审视了深度学习漏洞检测的评估方式，指出数据重复、跨项目划分不当和标签噪声会显著高估模型性能，并给出更严格的评估流程~\cite{chakraborty2021reveal}。LineVul 将预训练 Transformer 用于行级漏洞预测，并在多个数据集上报告了较强基线~\cite{fu2022linevul}。CodeBERT 通过编程语言与自然语言的双模态预训练提供通用代码表示~\cite{feng2020codebert}。Risse 等进一步指出，仅用函数长度、词频等表面特征就可能在多个漏洞数据集上取得高分，说明“高分”未必来自漏洞语义~\cite{risse2024limits,risse2025topscore}。这些工作构成了“训练型检测器”的主要参照；本文不微调模型，而是评估现成 LLM 在提示下的判定行为，因此二者不是同一训练范式，但在任务定义和指标上可以对照。

\subsection{大语言模型与提示策略}
Codex 展示了大规模代码预训练模型按自然语言指令生成程序的能力，也推动了对代码模型零样本行为的系统研究~\cite{chen2021codex}。思维链提示引导模型先产生中间推理再给出答案，已被证明能改善多步推理任务~\cite{wei2022cot}；自一致性方法通过多次采样并投票提高答案稳定性~\cite{wang2023selfconsistency}。在安全领域，Pearce 等考察了 LLM 的零样本漏洞修复，指出生成结果在功能正确性与安全性上都可能存在明显局限~\cite{pearce2023zeroshot}。这些结果说明，LLM 的安全判定与修复不能只看最终输出，还需要检查其依据是否与代码事实一致。

少样本提示与上下文增强是两类常见改进方向。前者通过固定示例向模型传达任务边界和输出格式；后者把调用关系、全局变量、类型声明等上下文放入提示，帮助模型判断跨函数的数据流与生命周期。二者都会增加输入长度和费用，因此需要在同一预算口径下与零样本基线比较，而不是只报告精度。

\subsection{漏洞检测基准与评估质量}
SecVulEval 面向 LLM 的上下文感知评估，提供函数级标签、CVE 与 CWE 元数据、语句级变更标注以及执行环境、外部函数、函数参数、全局变量和类型声明等上下文~\cite{ahmed2026secvuleval,secvulevaldata}。PrimeVul 强调数据质量、去重、污染控制以及使用漏洞函数与修复函数进行配对评估，指出仅靠单标签分类容易掩盖模型行为~\cite{ding2025primevul}。Risse 等对漏洞检测基准的构造效度提出质疑，指出标签噪声、不可判定样本和表面特征会扭曲结论~\cite{risse2024limits,risse2025topscore}。这些工作共同说明：评估集如何构造、哪些字段可见、样本之间是否相关，会直接影响结论。

现有基准多用于比较不同模型或不同训练方法，较少在固定基础模型下并列比较提示策略，并把独立检测、配对联合判别、修复生成、费用与误差类型放在同一套协议中报告。本文在课程项目范围内补足这一比较视角；由于只使用一个基础模型和两个规模有限的子集，本文不主张对全部 LLM 或全部漏洞类型作出普适结论。

\subsection{修复生成与补丁质量评估}
自动程序修复（automatic program repair，APR）通常区分“plausible patch”和“correct patch”：前者能通过现有测试，后者才真正修复缺陷并保持预期行为。Smith 等指出，用同一套测试既生成又验证补丁会系统性高估修复成功率，许多“修复”只是过拟合测试~\cite{smith2015overfitting}；Gazzola 等的综述把修复流程概括为故障定位、补丁生成和补丁验证，验证阶段的核心正是“用什么 oracle 判断补丁正确”~\cite{gazzola2019apr}。在 C/C++ 场景中，强 oracle 通常需要可编译工程、测试用例或 PoC；本项目的函数片段不具备这些条件。

因此，本文只把修复生成作为“结构代理”评估：检查输出是否为完整函数、是否与输入不同、是否与参考修复在规范化文本上一致。这些指标能发现格式失败、空输出、原样返回和明显偏离，但不能证明漏洞已被消除、行为保持不变或没有引入新风险。报告在摘要、结果和局限性中反复强调这一点，避免把“文本接近参考修复”写成“修复正确”。

\subsection{传统分析工具、Juliet 基准与分层评估}
静态分析器通过抽象解释、符号执行或模式匹配发现缺陷。Clang Static Analyzer 以路径敏感的过程内/过程间检查覆盖内存泄漏与 C++ 对象生命周期~\cite{clangsa}；Cppcheck 以轻量规则覆盖 \texttt{memleak}、\texttt{memleakOnRealloc} 等模式~\cite{cppcheck}。二者的召回与误报差异很大，且对不可编译片段会直接失败，因此报告必须同时给出 coverage 与 strict 口径，而不能把编译失败写成“无漏洞”。动态工具以可执行程序为输入：AddressSanitizer 通过影子内存检测越界与释放后使用~\cite{serebryany2012asan}，LeakSanitizer 在其基础上做泄漏检测~\cite{lsan}，Valgrind Memcheck 通过动态二进制插桩检测 definite/indirect leak~\cite{nethercote2007valgrind}。动态工具的高精确率依赖“泄漏在该执行路径上可观测”；没有覆盖率证据时，clean 只能读作“未观察到”，不能读作安全。

Juliet Test Suite 由 NIST SARD 发布，用模板生成同一缺陷的 bad/good 变体，覆盖 CWE 与多种数据流形态~\cite{boland2012juliet,sard}。它的优势是标签由宏结构直接给出、可编译可运行；局限是合成、单 CWE、模板化，且不同变体之间近重复程度高。本文的 Tier B 用 Juliet v1.3 CWE-401 构建 200 个冻结单元，并按近重复簇隔离 dev/test；这只支持“同一 CWE-401 用例上的工具行为”结论，不支持外推到真实项目。

分层评估并非本文首创：静态/动态/LLM 的工具属性不同，把它们池化会制造不可解释的平均分。本文的增量是把“统一字段 + 分层主终点 + 不池化”写成可执行契约，并在同一契约下同时报告 coverage、strict、失败分类和本地计算/API 两种成本。这样，负面结果（例如某方法 coverage 低、某层 F1 不可定义）与正面结果一样可审计。

\section{数据}
\label{sec:data}

\subsection{数据来源与构建}
数据源为项目根目录的 \path{SecVulEval.xlsx} 的 \texttt{train} 工作表，对应 SecVulEval 基准及其数据发布~\cite{ahmed2026secvuleval,secvulevaldata}。本地工作簿包含 25{,}440 条函数记录，其中 10{,}998 条标记为有漏洞、14{,}442 条标记为无漏洞。构建脚本 \path{build_datasets.py} 检查原始索引、标签、必要文本字段和序列化列表字段；由于 7 条记录的 \texttt{changed\_statements} 无法解析，实际保留 25{,}433 条可用记录，排除索引与原因记录在 JSON 元数据中。缺失或损坏的修复目标不能构成有效配对，去除首尾空白后函数正文相同的版本对也被排除；最终可用于抽样的候选配对为 10{,}004 对，其中因修复目标缺失或格式错误排除 975 对，因正文相同排除 15 对。采样使用局部随机数生成器 \texttt{random.Random(42)}。

两个 JSON 均保留每条记录的 16 个原始字段和完整函数正文，包括 \texttt{idx}、\texttt{project}、\texttt{commit\_id}、\texttt{filepath}、\texttt{func\_name}、\texttt{func\_body}、\texttt{is\_vulnerable}、\texttt{cve\_list}、\texttt{cwe\_list}、\texttt{changed\_lines}、\texttt{changed\_statements}、\texttt{fixed\_func\_idx} 和 \texttt{context} 等。\texttt{changed\_lines}、\texttt{changed\_statements}、\texttt{cve\_list}、\texttt{cwe\_list} 转为 JSON 数组；\texttt{context} 保留原始字符串，不执行其中的表达式。

\subsection{独立函数集}
独立函数集记为
\begin{equation}
\Dind=\{(x_i,y_i)\}_{i=1}^{200},\qquad y_i\in\{0,1\},
\end{equation}
其中 $x_i$ 为函数记录，$y_i=1$ 表示有漏洞。该集合严格包含 160 条有漏洞函数和 40 条无漏洞函数，并满足三条约束：
\begin{enumerate}
    \item 每个 \texttt{(project, commit\_id)} 分组最多选取一条函数记录；
    \item 任一入选漏洞函数在数据中标注的修复版本，不同时出现在该集合中；
    \item 40 条无漏洞函数均不是原始工作簿中任意漏洞记录通过 \texttt{fixed\_func\_idx} 指向的修复版本。
\end{enumerate}
这些约束保证集合内不存在数据集已标注的漏洞/修复配对关系，但不对上游完整 Git 历史作出证明。这里的“独立”指抽样与修复关联约束，不表示样本在统计意义上完全独立。

在 200 条源记录上冻结 \texttt{secvul-split-v1} 的检测划分：40 条开发（32 正/8 负）与 160 条测试（128 正/32 负）。划分按项目级互斥，40 个开发项目与 40 个测试项目没有交集；多例项目全部进入测试集。本文的模型结果只报告 160 条测试函数，40 条开发函数不进入主结果表。冻结算法、CWE 家族/长度分层和校验记录见 \path{splits/README.md} 与 \path{splits/split_v1.json}。

\subsection{漏洞/修复配对集}
配对集记为
\begin{equation}
\Dpair=\{(x_j^{v},x_j^{f})\}_{j=1}^{200},
\end{equation}
其中 $x_j^{v}$ 为标记有漏洞的版本，$x_j^{f}$ 为其对应修复且标记无漏洞的版本。配对必须满足
\begin{equation}
\texttt{fixed\_func\_idx}(x_j^{v})=\texttt{idx}(x_j^{f}),
\end{equation}
并且两个版本的项目、提交标识、文件路径和函数名一致，去除首尾空白后的函数正文不同。每个“项目—提交”分组最多选择一对。该集合包含 200 对、共 400 条函数记录。版本先后顺序由原始索引关联确定，不通过提交时间或文本相似度推测。

配对源集合冻结为 20 对开发 + 180 对测试，同一对的两个版本始终留在同一侧；20 个开发项目与 60 个测试项目在配对任务内没有交集。主结果使用 180 对测试版本，20 对开发版本只用于修复示例和 pilot。

\subsection{集合隔离与描述统计}
设 $I(\mathcal{D})$ 为集合中的原始记录索引，$G(\mathcal{D})$ 为“项目—提交”分组，则构建过程保证
\begin{align}
I(\Dind)\cap I(\Dpair)&=\varnothing,\\
G(\Dind)\cap G(\Dpair)&=\varnothing.
\end{align}
两个源集合合计包含 600 条不同的原始记录。冻结划分后，80 条不同记录进入开发（40 条检测 + 20 对配对），520 条进入测试（160 条检测 + 360 条配对函数）；开发与测试的源索引交集为 0，函数正文哈希交集为 0。检测任务内开发/测试项目交集为 0，配对任务内也为 0；但两个任务独立划分，因此有 2 个检测开发项目出现在配对测试（yara、znc），5 个配对开发项目出现在检测测试（Fast-DDS、OpenSC、\texttt{file}、htmldoc、libgd）。这不是完全的项目级隔离，这一点在外部有效性讨论中保留。按源 \texttt{project} 字段统计，两个源子集共享 26 个项目，项目并集为 134 个。

\begin{table}[t]
\caption{源子集与冻结划分的组成}
\label{tab:data-composition}
\centering
\small
\begin{tabular}{@{}lcrrr@{}}
\toprule
子集 & 评估单元 & 有漏洞 & 无漏洞 & 函数总数\\
\midrule
源独立集 & 200 条函数 & 160 & 40 & 200\\
源配对集 & 200 对版本 & 200 & 200 & 400\\
冻结检测开发集 & 40 条函数 & 32 & 8 & 40\\
冻结检测测试集 & 160 条函数 & 128 & 32 & 160\\
冻结配对开发集 & 20 对版本 & 20 & 20 & 40\\
冻结配对测试集 & 180 对版本 & 180 & 180 & 360\\
\bottomrule
\end{tabular}
\end{table}

表~\ref{tab:data-stats} 给出源子集的描述统计。函数行数按 \texttt{func\_body} 的换行符计数，包含空行；分布呈明显长尾，因此同时报告中位数、四分位区间和最大值，均值不作为主要概括。源独立集采用 160/40 的人为比例，源配对集因配对要求而标签平衡；两者的总体指标不能直接合并解释。

冻结测试集是源子集的子集：检测 160 条的函数行数中位数为 44.5（IQR 19.8--94.0，最大 1{,}172），配对漏洞版为 43.5（21.8--86.0，最大 1{,}011），修复版为 44.5（23.8--86.0，最大 1{,}015）。测试集的标签比例与源子集相同（128/32；180 对平衡），因此下文仍按 80/20 和配对平衡解释指标，但所有主结果的分母改为测试集。

\begin{table}[t]
\caption{源子集的描述统计}
\label{tab:data-stats}
\centering
\scriptsize
\begin{tabular}{@{}lrrr@{}}
\toprule
指标 & $\Dind$ & $\Dpair$ 漏洞 & $\Dpair$ 修复\\
\midrule
函数记录数 & 200 & 200 & 200\\
项目数 & 80 & 80 & 80\\
提交分组数 & 200 & 200 & 200\\
行数中位数 & 46 & 43.5 & 43.5\\
行数四分位区间 & 21.5--94.2 & 21.8--82.2 & 23.8--85.0\\
最大行数 & 1{,}172 & 1{,}011 & 1{,}015\\
含 CVE 编号记录 & 200 & 200 & 200\\
changed\_statements 非空 & 188 & 183 & 190\\
\bottomrule
\end{tabular}
\end{table}

\subsection{数据审计与泄漏风险}
只读审计发现若干会直接影响实验设计的事实，汇总在表~\ref{tab:audit}。这些事实不是模型结果，而是数据构造与字段分布的结果；报告在解释方法差异时必须把它们作为混淆因素。

\begin{table}[t]
\caption{数据审计中的关键事实}
\label{tab:audit}
\centering
\small
\begin{tabular}{@{}p{4.6cm}p{3.2cm}@{}}
\toprule
发现 & 影响\\
\midrule
5 条函数正文恰好 32{,}767 字符且在中途截断：源子集独立 3 条、配对 1 对；这 5 条全部落在测试集（独立 3 条 + 配对 1 对） & 长度分布是源截断后的下界；超长函数可能被截断 \\
源独立集 \texttt{context} 分层：非空 156/156 为漏洞，空 44 条中 4 条漏洞、40 条非漏洞；测试集非空 126/126 为漏洞，空 34 条中 2 条漏洞、32 条非漏洞 & context 非空层为单类别，与标签高度相关；M3 的提升可能来自格式差异而非语义 \\
源独立集 \texttt{fixed\_func\_idx} 非空：有漏洞 143/160，无漏洞 0/40；测试集为 115/128 与 0/32 & 存在性直接泄漏标签，必须从所有提示中排除 \\
源子集 CWE-401 漏洞函数 4/360（独立 2/160，配对 2/200）；测试集 3/308；严格 C++ 单标签 0 条 & 只能作案例系列，不能支持 CWE-401 或 Juliet 专项结论 \\
源独立集 CWE/CVE 非空：无漏洞 40/40（测试集 32/32）；配对修复版 200/200（测试集 180/180）复制漏洞版 & CWE/CVE 是溯源信息，不是函数级标签特征 \\
\bottomrule
\end{tabular}
\end{table}

第一项意味着样本长度存在源级截断：32{,}767 是 Excel 单元格字符上限，源子集 5 条记录均在函数中途结束，且这 5 条全部落在测试集。第二项是最重要的混淆：源独立集 context 非空层 156 条全部为漏洞样本，空层 44 条中 4 条漏洞、40 条非漏洞；冻结测试集非空层 126/126 为漏洞，空层 34 条中 2 条漏洞、32 条非漏洞。非空层是单类别，macro-F1 不可定义，M3 的总体差异与 context 可用性高度混杂。即使运行器只提取白名单字段，M3 仍然可以通过“有没有上下文”区分大量样本。因此本文对 M3 的结论只能限定为“在该数据上，上下文增强提示的整体效果”，不能声称模型理解了项目上下文。结果部分必须按 \texttt{context} 是否非空分层报告，至少给出该分层的样本数和指标；若 M3 的优势只出现在有上下文的样本上，也不能排除标签相关性的解释。第三项说明 \texttt{fixed\_func\_idx} 的存在性与标签相关，协议一律排除该字段；第四项说明 CWE/CVE 不能作为输入特征，配对修复版复制漏洞版的 CWE/CVE 也意味着不能用它区分对内的两个版本。

\subsection{研究范围：全 CWE 定位与 CWE-401 样本量}
\label{sec:scope-cwe}
表~\ref{tab:audit} 已给出 CWE-401 的总量，本节说明计数口径与它对研究定位的影响。\texttt{cwe\_list} 中存在 \texttt{CWE-415CWE-401} 这类连写标签，统计时按正则 \texttt{CWE-[\_ ]?(\textbackslash d+)} 拆分，并在单条记录内去重；只有 \texttt{is\_vulnerable=true} 的记录才计入漏洞样本，修复版本和独立集负例即使保留 CWE-401 元数据也不计入。按此口径，源独立集有 2 条 CWE-401 漏洞函数（\texttt{idx}=17743、5896，占 2/160），源配对集有 2 条（\texttt{idx}=5886、8675，占 2/200），合计 4/360（1.11\%）。冻结测试集有 3 条（5896、5886、8675），占 3/308；17743 落在检测开发集。4 条源记录中 3 条是 \texttt{.c}，1 条是 \texttt{.cc}；唯一的 C++ 记录 17743 的原始标签是 \texttt{CWE-415CWE-401}，因此“C++ 扩展名且单一 CWE-401 标签”为 0 条。4 条记录只来自 3 个项目（linux 2 条，dcmtk 1 条，ImageMagick 1 条），其中 8675 的函数正文恰好 32{,}767 字符并在中途截断。

CWE 的整体分布同样不支持逐类推断。两个源子集的漏洞函数合并后只有 57 种数字 CWE；源独立集 39 种中有 23 种只出现 1 次，源配对集 39 种中有 18 种只出现 1 次；源独立集只有 6 种、源配对集只有 7 种 CWE 达到 $n\ge 10$。冻结测试集合并后为 54 种，其中 29 种只出现 1 次，只有 10 种达到 $n\ge 10$。CWE-401 在测试集的合并计数为 3，仍属于长尾类别。若对每个 CWE 报告 precision/recall，大多数类别的分母只有 1--9 条，一条预测翻转就会主导结果。

因此本文的主定位是跨 CWE 的函数级评估，而不是 CWE-401、C++ 内存泄漏或 Juliet CWE-401 专项。源子集的 4 条 CWE-401 记录（测试集 3 条）只作为案例系列：用于观察模型是否提到错误路径上的资源释放和生命周期，不进入逐 CWE 排名，不估计 CWE-401 的召回率或 F1，也不与 Juliet 结果比较。数据中大小写不敏感搜索 \texttt{juliet} 的命中数为 0，600/600 条记录都带 CVE，进一步说明当前子集是真实项目 CVE 抽样，不是合成弱点测试集。

\subsection{开发示例、输入字段与任务矩阵}
三种策略中的少样本提示需要示例。冻结划分 \texttt{secvul-split-v1} 从检测开发集取 4 条示例（2 条有漏洞、2 条无漏洞），从配对开发集取 2 条修复示例（每条含漏洞版本与参考修复），保存在 \path{splits/dev_examples_fewshot.json}；示例的代码、标签、理由和来源索引由 \path{splits/example_notes.json} 冻结。示例只来自开发集，不来自 160 条检测测试函数或 180 对测试版本。运行器加载时对评估文件执行隔离检查：示例的规范化源码不得与任何测试函数完全相同，示例声明的 \texttt{source\_idx} 不得出现在测试集中；检测示例必须覆盖正负两类，修复示例必须包含完整修复。示例的标签和参考输出是方法的一部分，不是测试样本的标签泄漏。

开发集与测试集按项目级互斥构造：检测 40 个开发项目与 40 个测试项目没有交集，配对 20 个开发项目与 60 个测试项目没有交集；开发/测试源索引交集为 0，函数正文哈希交集为 0。两个任务独立划分，因此有 2 个检测开发项目出现在配对测试、5 个配对开发项目出现在检测测试；这不会让示例代码进入同任务的测试集，但跨任务并非完全项目隔离。冻结记录、源文件 SHA-256 和各 artifact 哈希见 \path{splits/FREEZE.txt} 与 \path{splits/split_v1.json}。

课程指南建议 40 条开发样本加 160 条测试样本，冻结划分与此一致（检测 40/160；配对任务按对计数，20/180）。开发集只用于固定示例、pilot 和人工校准，不用于报告主结果；测试集在提示和示例冻结后才可运行。示例不按测试表现挑选，提示版本和示例哈希写入运行 manifest。若需要改划分或示例，必须生成新的 \texttt{split\_id} 和新的 freeze 记录，不能覆盖旧文件后重跑。

M1 与 M2 的目标函数输入只包含 \texttt{func\_body}；M3 另加第~\ref{sec:methods}节定义的、按白名单选取的上下文。以下字段不得进入任何策略的提示：\texttt{is\_vulnerable}（标签）、\texttt{cve\_list} 与 \texttt{cwe\_list}（漏洞类型元数据）、\texttt{commit\_message}（可能描述修复）、\texttt{changed\_lines} 与 \texttt{changed\_statements}（直接指出变更位置）、\texttt{fixed\_func\_idx} 与 \texttt{idx}（可用于反查标签）、\texttt{project} 与 \texttt{filepath}（可用于检索公开修复），以及 \texttt{context} 中的 \texttt{Explanation} 字段（可能包含漏洞说明）。本地数据中该字段当前均为空，但协议仍将其排除，以避免数据更新后产生泄漏。

配对集的评分需要两个版本，但推理时每个版本独立调用模型，模型看不到另一个版本，也不知道当前函数属于哪一对。配对关系只在评分阶段使用。这样，配对指标衡量的是模型对同一函数两种安全状态的联合判别，而不是模型对“找不同”任务的直接比较能力；它也避免了把两个版本放在同一次调用中可能引入的位置偏差。

表~\ref{tab:field-visibility} 汇总各字段在三种策略中的可见性。该表是防止标签泄漏的主要依据：任何标为“否”的字段都不得以文本、元数据或文件名形式进入提示。若后续实验需要加入新字段，必须先更新该表并在报告中说明理由，不能在没有记录的情况下临时追加。

\begin{table*}[t]
\caption{原始字段在三种策略提示中的可见性}
\label{tab:field-visibility}
\centering
\small
\begin{tabular}{@{}>{\raggedright\arraybackslash}p{4.8cm}lll>{\raggedright\arraybackslash}p{6.0cm}@{}}
\toprule
字段 & M1 & M2 & M3 & 处理说明\\
\midrule
\texttt{func\_body} & 是 & 是 & 是 & 唯一必选输入，保留完整函数正文\\
\texttt{context} 白名单 & 否 & 否 & 是 & 仅执行环境、外部函数、函数参数、全局变量、类型声明\\
\texttt{context.Explanation} & 否 & 否 & 否 & 可能包含漏洞说明，协议上一律排除\\
\texttt{is\_vulnerable} & 否 & 否 & 否 & 预测目标，不得泄露\\
\texttt{cve\_list} / \texttt{cwe\_list} & 否 & 否 & 否 & 漏洞类型元数据，可能直接提示答案\\
\texttt{commit\_message} & 否 & 否 & 否 & 可能描述修复内容或安全影响\\
\texttt{changed\_lines} / \texttt{changed\_statements} & 否 & 否 & 否 & 直接指出修复位置，构成强泄漏\\
\texttt{idx} / \texttt{fixed\_func\_idx} & 否 & 否 & 否 & 可用于反查原始记录与标签\\
\texttt{project} / \texttt{filepath} / \texttt{func\_name} & 否 & 否 & 否 & 可用于检索公开修复或根据命名猜测\\
\bottomrule
\end{tabular}
\end{table*}

表~\ref{tab:task-matrix} 给出两个任务在冻结测试集上的请求数。检测任务对 160 条独立测试函数和 360 条配对测试函数各请求一次；修复任务只对漏洞版本请求，独立测试集 128 条、配对测试集 180 条。三种策略各运行一遍，每个策略 828 次请求，三种策略合计 2{,}484 次请求。该数字是规划值，不是实测调用数；实际调用数、重试和缓存命中以运行日志为准。

\begin{table}[t]
\caption{任务与请求矩阵（每种策略，冻结测试集）}
\label{tab:task-matrix}
\centering
\small
\begin{tabular}{@{}lrrr@{}}
\toprule
数据集 & 检测请求 & 修复请求 & 评估单元\\
\midrule
独立测试集 & 160 & 128 & 160 条函数\\
配对测试集 & 360 & 180 & 180 对 / 360 条函数\\
\midrule
合计 & 520 & 308 & 160 条函数 + 180 对\\
\bottomrule
\end{tabular}
\end{table}

\section{方法}
\label{sec:methods}

\subsection{任务形式化与输出模式}
检测任务要求模型对一条函数 $x$ 输出二值判断 $\hat{y}(x)\in\{0,1\}$，以及可选的置信度和简短理由。修复任务要求模型对一条漏洞函数 $x^{v}$ 输出一个完整的修复函数 $\hat{x}^{f}$ 和解释。两个任务使用同一基础模型和同一提示骨架，但输出模式不同。

检测输出必须是可解析的 JSON 对象，包含且仅要求以下字段：
\begin{lstlisting}
{"is_vulnerable": true|false,
 "confidence": 0.0-1.0,
 "rationale": "brief reason"}
\end{lstlisting}
修复输出必须是：
\begin{lstlisting}
{"fixed_function": "<complete function>",
 "explanation": "brief explanation"}
\end{lstlisting}
其中 \texttt{fixed\_function} 必须包含完整函数定义、签名和花括号，不能是 diff、片段或 Markdown。评分只读取 \texttt{is\_vulnerable} 和 \texttt{fixed\_function}；\texttt{confidence}、\texttt{rationale} 和 \texttt{explanation} 用于定性分析，不参与主指标。把置信度排除在主指标之外，是因为不同模型对概率标定的含义不一致，而且强制输出概率会诱导过度自信。

对配对 $(x_j^{v},x_j^{f})$，检测任务分别得到 $\hat{y}(x_j^{v})$ 和 $\hat{y}(x_j^{f})$。理想结果是 $(\hat{y}^{v},\hat{y}^{f})=(1,0)$。若模型对两个版本都输出 0，它在漏洞版本上犯错；若都输出 1，它在修复版本上犯错；若输出 $(0,1)$，则是“反向翻转”。配对关系只在评分阶段使用，推理时两个版本各自独立调用，模型看不到另一个版本。

\subsection{任务边界}
本文评估的是函数级、单标签的静态判定与修复生成。检测不要求定位到具体行；修复不执行代码、不编译、不运行测试，也不进行跨文件构建、指针分析或污点追踪。因此，本文不能回答“漏洞在哪一行”“补丁是否真正消除漏洞”“补丁是否保持行为”等问题；这些问题需要行级标注、可编译工程和测试预言，超出当前数据与课程项目的范围。

这一边界也影响指标解释。函数级标签是二值的：只要函数中存在一个数据集标注的缺陷，标签即为 1；模型不需要找出全部缺陷，也不因漏掉第二个缺陷而额外扣分。配对集的修复版本标签为 0，但修复版本可能仍含其他未标注问题。修复生成的结构指标只能说明输出“像不像一个完整补丁”，不能说明“补丁是否正确”。报告在解释假阳性、假阴性和修复结果时，必须把“标签定义”“结构代理”与“代码实际风险”区分开。

\subsection{共享控制}
三种策略刻意保持以下条件一致，以隔离提示策略的影响：
\begin{itemize}
    \item \textbf{基础模型相同}：模型名称、版本或快照、服务商和访问日期记录在表~\ref{tab:model}；三种策略在同一时间段内运行，若服务端模型发生更新则整批重跑。
    \item \textbf{解码参数相同}：温度固定为 0，\texttt{top\_p} 为 1，最大输出 token 按任务固定；主实验关闭思考模式（若端点支持）。不启用工具调用、联网或代码执行。
    \item \textbf{调用粒度相同}：每次只提交一条函数；配对集的两个版本分两次调用，不把同一对放在同一次请求中。
    \item \textbf{输出模式相同}：检测和修复分别使用固定 JSON 模式；三种策略的字段要求一致。
    \item \textbf{解析与重试相同}：传输错误、429 和 5xx 按指数退避重试；HTTP 200 但输出不符合模式时记为 invalid，不自动重问。续跑时可选择重试 invalid，但该选择必须记录。
    \item \textbf{提示语言相同}：三种策略的任务说明、示例和输出要求均使用英文，避免语言切换带来的额外差异。
\end{itemize}
不强制三种提示等长。少样本示例和上下文必然增加输入 token，这部分成本差异本身就是比较对象；若强行截断到相同长度，反而会改变策略定义。报告时分别给出原始 token 数与费用。

\subsection{策略 M1：零样本直接判定}
M1 只提供任务说明和函数源码。检测提示要求模型直接给出 JSON 结论，不提供示例，不提供调用关系或全局变量等上下文；修复提示要求模型返回完整修复函数。M1 是课程指南中“直接提示”的对应实现，也是另外两种策略的增量基线。M1/M2 的检测提示模板见附录~\ref{app:prompts}。

\subsection{策略 M2：固定少样本提示}
M2 在 M1 的任务说明之后加入固定示例，再给出待判定函数。检测任务使用 4 条开发集示例（2 条有漏洞、2 条无漏洞），修复任务使用 2 条开发集配对示例（每条含漏洞版本和参考修复）。示例来自 \path{splits/dev_examples_fewshot.json}，由 \path{splits/example_notes.json} 冻结，选取规则在实验前确定：
\begin{enumerate}
    \item 示例只从 40 条检测开发集和 20 对配对开发集中选取，不来自 160 条检测测试或 180 对测试版本；
    \item 运行器检查示例的规范化源码不与任何测试函数相同，示例声明的 \texttt{source\_idx} 不在测试集中；
    \item 检测示例覆盖正负两类，修复示例包含完整函数和参考修复；
    \item 示例顺序固定，不在每次调用时重新抽样。
\end{enumerate}
示例的标签和输出是方法的一部分；它们来自开发集，不是测试样本的标签泄漏。M2 与 M1 的差异仅在于示例，因此二者的差值可解释为“少样本示范”的效果。

\subsection{策略 M3：白名单上下文增强}
M3 在 M1 的输入上增加 SecVulEval 的上下文白名单字段：\texttt{Execution Environment}、\texttt{External Function}、\texttt{Function Argument}、\texttt{Globals} 和 \texttt{Type Execution Declaration}。原始 \texttt{context} 是 Python 字典的字符串表示，其中数组以 \texttt{array([...], dtype=object)} 形式出现，不能直接当作 JSON 解析。运行器按字段名提取白名单内容，逐项输出为 \texttt{- item} 列表，只保留非空字段；若五个字段全为空，则输出 \texttt{(No usable structured context was available.)}，此时 M3 在该样本上退化为 M1；若整体解析失败，则该样本不附带上下文，并在日志中计数。

M3 不使用 \texttt{Explanation}、CVE/CWE 标签、变更标注或提交信息，也不调用 Clang 等外部静态分析工具。这样选择的原因是：评估样本是来自不同项目的函数片段，很多片段无法独立编译，外部工具会产生大量与样本无关的缺失依赖告警，反而引入噪声；此外，课程指南的 L3 示例使用 Clang 分析结果，但这些函数片段通常缺少头文件、宏和类型定义，难以作为独立翻译单元分析，因此改用数据自带的项目上下文。该替换在报告中明确说明，避免把不同实现当作同一方法。

M3 有一个无法通过提示清洗消除的混淆：在冻结检测测试集上，context 非空层 126/126 为漏洞样本，空层 34 条中 2 条漏洞、32 条非漏洞（源独立集非空 156/156、空 44 条中 4 条漏洞）；“上下文是否可用”与标签高度相关，且非空层是单类别。配对测试集内两个版本的 context 相同（176 对两侧非空、4 对两侧为空），该混淆较弱，但配对集的结果仍不能说明模型理解了上下文。就检测测试集而言，M3 相对于 M1 的提升可能来自“有没有上下文”这一格式差异，而不是模型理解了项目上下文。结果部分必须按 context 是否非空分层报告，至少给出分层样本数和指标；M3 的主结论只能限定为“在该数据上，上下文增强提示的整体效果”。

\subsection{与课程指南三种策略的对应关系}
课程指南给出的示例策略为 L1 直接提示、L2 少样本提示、L3 分析引导提示，并要求在同一批测试样本上运行。本文的 M1、M2、M3 分别对应这三类策略，但把 L3 的“Clang 分析结果”替换为“SecVulEval 白名单上下文”，原因如上。三种策略共享同一基础模型，属于“比较提示策略”的设计，而不是“比较模型”的设计；若后续增加第二个基础模型，应作为独立实验报告，不与当前结果混合。

\section{实验设置}
\label{sec:eval}

\subsection{评估规模与运行}
按表~\ref{tab:task-matrix}，每个策略共 520 次检测请求和 308 次修复请求，合计 828 次；三种策略合计 2{,}484 次，六次运行均已完成。正式运行前曾在 \path{splits/pilot_detection_36.json}（36 条检测）和 \path{splits/pilot_pairs_18.json}（18 对）上做端到端试运行，检查提示渲染、上下文解析、JSON 解析、重试和计费字段；pilot 不进入结果表。运行中传输错误按指数退避重试并单独计数，不覆盖原始失败记录；输出不符合模式时记为 invalid，不通过反复重问掩盖失败率。

\subsection{模型与解码配置}
表~\ref{tab:model} 记录实际使用的基础模型与调用配置。六次正式运行（三策略 × 两数据集）的 manifest 与请求体一致：模型 ID 为 \texttt{deepseek-v4.1-flash}，服务端为 OpenCode Go 端点，响应中的 \texttt{model} 为同一 ID，\texttt{system\_fingerprint} 为 \texttt{null}（服务端未返回快照标识）。请求体明确写入 \texttt{"thinking": \{"type": "disabled"\}}，温度为 0，最大输出 8192 token，并发 8；\texttt{top\_p} 未显式设置，按服务端默认行为处理，不能假定为某个固定值。价格快照为输入 0.30、输出 1.20、缓存输入 0.006 美元每百万 token；所有运行在 2026-11-11（UTC）完成，运行状态均为 completed。

\begin{table}[t]
\caption{基础模型与调用配置}
\label{tab:model}
\centering
\small
\begin{tabular}{@{}>{\raggedright\arraybackslash}p{3.0cm}>{\raggedright\arraybackslash}p{5.2cm}@{}}
\toprule
项目 & 配置\\
\midrule
基础模型 & \texttt{deepseek-v4.1-flash}\\
API ID & \texttt{deepseek-v4.1-flash}\\
服务商 / 端点 & OpenCode Go；\path{https://opencode.ai/zen/go/v1/chat/completions}\\
响应版本/快照 & \texttt{model}=\texttt{deepseek-v4.1-flash}；\texttt{system\_fingerprint}=null\\
温度 / \texttt{top\_p} & 0 / 未显式设置（服务端默认）\\
思考模式 & disabled（请求体显式设置）\\
最大输出 token & 8{,}192\\
并发度 & 8\\
输出格式 & JSON 对象（\texttt{json\_object}）\\
每次调用函数数 & 1\\
重试策略 & 传输错误最多 3 次；schema-invalid 不自动重问\\
价格快照 & 输入 0.30 / 输出 1.20 / 缓存输入 0.006 USD 每 1M token\\
运行日期 & 2026-11-11（UTC）\\
\bottomrule
\end{tabular}
\end{table}

\subsection{检测指标}
对有漏洞类（正类）定义 TP、FP、TN、FN，使用精确率 $P$、召回率 $R$、F1、准确率、特异度、平衡准确率和 Matthews 相关系数（MCC）：
\begin{align}
P&=\frac{TP}{TP+FP}, \qquad R=\frac{TP}{TP+FN},\\
F_1&=\frac{2PR}{P+R},\\
\mathrm{Acc}&=\frac{TP+TN}{TP+TN+FP+FN},\\
\mathrm{Spec}&=\frac{TN}{TN+FP},\\
\mathrm{BalAcc}&=\frac{R+\mathrm{Spec}}{2},\\
\mathrm{MCC}&=\frac{TP\cdot TN-FP\cdot FN}{D},\\
D&=\sqrt{(TP+FP)(TP+FN)}\notag\\
&\quad\times\sqrt{(TN+FP)(TN+FN)}.
\end{align}
当分母为 0 时，相应指标记为未定义（报告为 ---），MCC 的分母为 0 时记 0。独立集的有漏洞比例为 80\%，准确率会被多数类主导；“全判有漏洞”即可得到 80\% 准确率，但特异度为 0、MCC 为 0。因此以 macro-F1 和 MCC 为主要指标，精确率、召回率、特异度和平衡准确率分别报告。

macro-F1 定义为两个类别 F1 的算术平均：
\begin{equation}
\mathrm{MacroF1}=\frac{1}{2}\left(F_1^{\mathrm{vuln}}+F_1^{\mathrm{non}}\right),
\end{equation}
其中 $F_1^{\mathrm{vuln}}$ 以有漏洞类为正类，$F_1^{\mathrm{non}}$ 以无漏洞类为正类。该指标由混淆矩阵计算，对类别不平衡不敏感。

检测输出有两种分析口径。\textbf{valid-only} 只统计通过模式校验的预测，用于比较“模型在能给出有效输出时的判别能力”；\textbf{strict all-samples} 把所有 invalid 和 API 错误都当作错误，分母为全部计划请求，用于比较“端到端可用性”。每个结果表必须同时给出有效预测数、invalid 数、API 错误数和覆盖率；若某策略的 invalid rate 明显更高，valid-only 指标会虚高，不能只报 valid-only。

\subsection{配对联合指标}
配对集的检测请求按对聚合。对每一对 $j$，根据两个版本的预测定义四类联合结果：
\begin{itemize}
    \item \textbf{正确区分}（\texttt{vulnerable\_detected\_fixed\_clean}）：漏洞侧判 1、修复侧判 0；
    \item \textbf{双报漏洞}（\texttt{both\_flagged}）：两侧都判 1；
    \item \textbf{双报安全}（\texttt{both\_clean}）：两侧都判 0；
    \item \textbf{漏报加误报}（\path{vulnerable_missed_fixed_flagged}）：漏洞侧判 0、修复侧判 1。
\end{itemize}
联合判别准确率定义为正确区分对数除以配对总数：
\begin{equation}
\mathrm{DA}=\frac{N_{\mathrm{correct}}}{N_{\mathrm{pair}}},
\end{equation}
其中 $N_{\mathrm{correct}}$ 为漏洞侧判 1、修复侧判 0 的对数，$N_{\mathrm{pair}}$ 为参与评分的对数。同时报告 valid-pairs 版本（只在两侧都有效的对上计算）和 all-pairs 版本（把无效对当作失败）。四类结果的比例以有效对为分母，并给出计数。仅报告版本级准确率会掩盖配对内部的相关性；一个总是判“无漏洞”的模型可以在修复侧得到 100\% 正确，但漏洞侧全部漏报，联合判别准确率为 0。

\subsection{修复结构指标}
修复任务只对漏洞版本请求。由于没有编译、测试或 PoC oracle，本文只报告以下结构代理：
\begin{itemize}
    \item \textbf{schema 有效率}：输出可解析且包含 \texttt{fixed\_function} 的比例；
    \item \textbf{完整函数率}：在 schema 有效输出中，括号配对、存在函数式前缀、不是 diff 或 Markdown 的比例；
    \item \textbf{与输入不同率}：在 schema 有效输出中，规范化后与输入函数不同的比例；
    \item \textbf{完整且不同率}：同时满足完整函数和与输入不同的比例；
    \item \textbf{精确匹配参考修复率}：在配对集有参考修复的漏洞版本上，规范化文本与参考修复完全相同的比例；
    \item \textbf{空白折叠匹配率}：折叠空白后与参考修复相同的比例。
\end{itemize}
这些指标的分母不同，报告时必须写清：完整函数率、与输入不同率、完整且不同率以 schema 有效输出为分母；精确匹配和空白折叠匹配只对配对测试集的 180 条漏洞版本有参考修复，独立测试集的 128 条漏洞版本没有参考修复。所有指标都是结构代理，不能解释为“修复正确率”或“漏洞消除率”。

\subsection{成本与延迟}
设第 $i$ 次调用消耗 $t_i^{\mathrm{in}}$ 个输入 token 和 $t_i^{\mathrm{out}}$ 个输出 token，服务商单价分别为 $p_{\mathrm{in}}$ 和 $p_{\mathrm{out}}$（美元/百万 token），则总费用为
\begin{equation}
C=\sum_{i=1}^{n}\left(t_i^{\mathrm{in}}p_{\mathrm{in}}+t_i^{\mathrm{out}}p_{\mathrm{out}}\right).
\end{equation}
若服务商报告缓存命中 token，则分别按缓存命中和未命中单价计算，并同时报告“全按未命中价”的保守上界。单价按运行当日服务商公开价格记录，并注明币种、是否含税和是否有折扣。重试调用计入总费用，但不计入主结果的样本数；开发示例调用单独汇总。若未配置价格，运行器的成本字段为 \texttt{null}，报告不得用估算值替代实测账单。

成本—效果指标包括：
\begin{equation}
C_{100}=\frac{C}{N_{\mathrm{unit}}}\times 100,\qquad
C_{\mathrm{correct}}=\frac{C}{N_{\mathrm{correct}}},
\end{equation}
其中 $N_{\mathrm{correct}}$ 在检测任务中为 $TP+TN$，在配对任务中为正确区分对数，在修复任务中为“完整且不同”的样本数。若分母为 0，则相应成本记为不适用（---）。延迟报告 p50 和 p95，并注明并发度、测量地点和是否包含重试等待。

\subsection{统计推断}
对 accuracy、recall、specificity、联合判别准确率、schema 有效率等比例，使用 95\% Wilson 区间；对 F1、macro-F1、MCC 和策略差值，使用 bootstrap 区间。检测任务按项目聚类重采样，配对任务按对重采样，修复任务按对或按函数重采样；同一对的两个版本始终一起重采样。重复 10{,}000 次，随机种子固定为 42，报告中位数和 2.5/97.5 百分位。

策略间的成对比较使用 McNemar 精确检验：对每条独立集记录判断两种策略是否都判对，仅统计一方正确、另一方错误的 discordant 样本；对配对集，以每一对的联合判别是否正确为二元结果；对修复任务，以“完整且不同”是否成立为二元结果。三种策略产生 3 组两两比较，在每个任务内用 Holm--Bonferroni 控制多重比较。McNemar 假设单元独立，同项目单元可能相关，因此把 McNemar 作为次要检验，主推断使用按项目或按对聚类的 bootstrap 差值区间。检验结果与点估计一起解读；若样本量不足以支持显著性，明确说明只描述观测差异，不宣称优劣。

\subsection{实现与日志}
评估脚本、提示模板、解析器和评分脚本固定版本，并与结果表一一对应。运行器 \path{runner/run_experiment.py} 通过 OpenAI 兼容的 Chat Completions 接口调用模型，只依赖 Python 标准库和 \texttt{httpx}；所有调用通过环境变量读取密钥，仓库中不保存密钥、请求头或完整服务商响应中的账户信息。逐样本结果保存为 JSON Lines，每行包含样本 ID、任务、策略、原始输出、解析结果、token 数、延迟和缓存信息；聚合脚本从该文件生成表格，禁止手工抄写数值。六次正式运行 \path{runner/runs/test_v1_{direct,few-shot,context}_{independent,pairs}} 已于 2026-11-11 完成，每策略 828 次请求、三策略合计 2{,}484 次；运行器版本 1.0.0，提示版本 \texttt{secvul-prompt-2026-11-11.1}。

表~\ref{tab:ind-results} 至表~\ref{tab:repair} 的行由 \path{report/scripts/fill_tables.py} 从六次运行的 \path{metrics.json} 生成，输出保存在 \path{report/generated_tables.txt}；检测/配对的区间、McNemar/Holm 和 context 分层由 \path{analysis/analyze_runs.py} 聚合为 \path{analysis/aggregation.json} 与 \path{analysis/aggregation.md}（10{,}000 次 project/pair cluster bootstrap，seed 42）。成本表由 \path{results.jsonl} 的记录级 \texttt{task} 字段汇总，输出在 \path{report/generated_costs.json}，而不是从 \path{metrics.json} 的跨任务 usage 拆分：每条记录读取 prompt/completion/cached token、\texttt{response.cost.cost\_usd}、\texttt{response.usage.opencode.allowance} 和 \texttt{response.request\_latency\_seconds}；\texttt{allowance} 是 OpenCode Go 套餐额度，不是 USD 成本，单独报告。只有全部计划记录都存在且每条记录都有数值成本时才给出总成本；本数据六个 run 都完整，因此成本有确定总值。partial 的 dev\_pilot 运行被排除在正式表之外。



\section{结果}
\label{sec:results}

\noindent\textit{说明：本节报告冻结测试集上的完整运行结果；参考基线为规则基线，不是 LLM。所有区间为 95\% 置信区间或 10{,}000 次 cluster bootstrap 区间。}

\subsection{参考基线：规则基线，不是 LLM}
\label{sec:ref-baselines}
在报告模型结果之前，本节先给出冻结测试集上由标签和构造规则直接计算的参考基线。它们来自 \path{baseline/results/test_split_baselines.json}，固定种子 42；脚本不调用模型、不访问网络、不使用凭证。这些行是\textbf{规则基线，不是 LLM 检测结果}，不能填入表~\ref{tab:ind-results} 至表~\ref{tab:repair} 的 M1--M3 行，也不参与模型间检验。它们只用于解释指标下限和配对机会水平。

\begin{table*}[t]
\caption{独立测试集上的参考基线（规则基线，非 LLM；160 条：128 正 / 32 负）}
\label{tab:ref-independent}
\centering
\small
\begin{tabular}{@{}lccccccccc@{}}
\toprule
基线 & TP & FP & FN & TN & Acc (\%) & P & R & F1 & Spec\\
\midrule
永远预测有漏洞（多数类） & 128 & 32 & 0 & 0 & 80.0 & 0.800 & 1.000 & 0.889 & 0.000\\
永远预测无漏洞 & 0 & 0 & 128 & 32 & 20.0 & --- & 0.000 & 0.000 & 1.000\\
\bottomrule
\end{tabular}
\end{table*}

表~\ref{tab:ref-independent} 的多数类基线等于测试集的有漏洞比例 128/160。它得到 80.0\% 准确率（Wilson 95\% 区间 73.1\%--85.5\%），但特异度为 0、平衡准确率为 0.5；按第~\ref{sec:eval}节的定义，MCC 的分母为 0 时记 0，因此该基线 MCC 为 0。永远预测无漏洞的准确率为 20.0\%（区间 14.5\%--26.9\%），精确率因没有正类预测而未定义，F1 在基线脚本中记为 0.000。两个常量基线说明：在 80/20 的测试集上，单看准确率无法区分“检测”与“多数类猜测”；任何模型结果都必须与特异度、平衡准确率、MCC 和 macro-F1 一起报告。

\begin{table*}[t]
\caption{配对测试集上的参考基线（规则基线，非 LLM；180 对 / 360 条函数记录）}
\label{tab:ref-paired}
\centering
\small
\begin{tabular}{@{}lcccccc@{}}
\toprule
基线 & Acc (\%) & P & R & F1 & 双方均正确 & 至少一方正确\\
\midrule
永远预测有漏洞（两侧） & 50.0 & 0.500 & 1.000 & 0.667 & 0/180 (0.0\%) & 180/180 (100.0\%)\\
永远预测无漏洞（两侧） & 50.0 & --- & 0.000 & 0.000 & 0/180 (0.0\%) & 180/180 (100.0\%)\\
配对内随机指定一侧为漏洞 & 56.1 & 0.561 & 0.561 & 0.561 & 101/180 (56.1\%) & 101/180 (56.1\%)\\
\bottomrule
\end{tabular}
\end{table*}

表~\ref{tab:ref-paired} 的“双方均正确”是配对联合判别的严格指标：只有漏洞侧判 1、修复侧判 0 才计入。两个常量基线在 360 条函数上的准确率都是 50.0\%，但双方均正确都是 0/180：每一对总有一个成员被固定判错。配对内随机指定一侧为漏洞得到 101/180（56.1\%，Wilson 95\% 区间 48.8\%--63.2\%），这是固定种子 42 的一次实现，解析期望为 50\%；它才是配对场景的机会水平。表中“至少一方正确”对常量基线为 100\%，进一步说明只看宽松指标会掩盖配对失败。

这些基线只使用标签和构造规则，不使用函数代码，因此不代表任何检测能力。它们的作用是给出解释锚点：测试集准确率的平凡下限是 80.0\%，配对联合判别的机会水平为 56.1\%（解析期望 50\%），常量基线的双方均正确率为 0。模型结果必须相对于这些锚点解释，但基线本身不进入模型排名、不计算 $p$ 值，也不与 M1--M3 合并。

\subsection{独立测试集上的函数级检测}
表~\ref{tab:ind-results} 汇总三种策略在独立测试集上的精确率、召回率、F1、准确率、特异度、平衡准确率、MCC、macro-F1 和覆盖率。三种策略都使用相同的 160 条测试记录和相同的评分脚本。macro-F1 由混淆矩阵计算；“全判有漏洞”可得到 80\% 准确率，但特异度和 MCC 均为 0，因此准确率不能单独作为结论。覆盖率均为 100\%，即所有记录都给出了可解析输出。few-shot 的 macro-F1 0.399 在数值上最高，但三种策略的区间高度重叠，不能据此排名。

\begin{table*}[t]
\caption{独立测试集上的检测结果（160 条：128 正 / 32 负）}
\label{tab:ind-results}
\centering
\scriptsize
\begin{tabular}{@{}lccccccccc@{}}
\toprule
策略 & P & R & F1 & Acc & Spec & BalAcc & MCC & Macro-F1 & Cov.\\
\midrule
M1 零样本直接判定 & 0.826 & 0.297 & 0.437 & 38.8 & 75.0 & 52.3 & 0.041 & 0.383 & 100.0\\
M2 固定少样本 & 0.837 & 0.320 & 0.463 & 40.6 & 75.0 & 53.5 & 0.061 & 0.399 & 100.0\\
M3 上下文增强 & 0.812 & 0.305 & 0.443 & 38.8 & 71.9 & 51.2 & 0.020 & 0.381 & 100.0\\
\bottomrule
\end{tabular}
\end{table*}

表中 P、R、F1、Acc、Spec、BalAcc、Cov. 分别为精确率、召回率、F1、准确率、特异度、平衡准确率和覆盖率；后续表格沿用这些缩写。对 32 条无漏洞测试函数，每一条误报都会使特异度变化约 3.1 个百分点，因此特异度的区间很宽。三种策略的 macro-F1 分别为 0.383（95\% project-cluster bootstrap 区间 0.305--0.552）、0.399（0.305--0.593）和 0.381（0.286--0.560）；MCC 也很低且区间跨 0。M3 的 context 分层结果见表~\ref{tab:m3-strata}：非空层 126 条全部为漏洞，macro-F1 不可定义；空层只有 34 条（2 漏洞/32 非漏洞），M1、M2、M3 的 macro-F1 分别为 0.512、0.512、0.494，M3−M1 的差值为 $-0.0179$（bootstrap 区间 $-0.1071$--0.0000）。因此 M3 与 M1 的总体差异不能解释为“上下文推理能力”。

\begin{table}[t]
\caption{独立测试集按 context 是否非空分层（探索性）}
\label{tab:m3-strata}
\centering
\scriptsize
\begin{tabular}{@{}lcccc@{}}
\toprule
分层 & n (V/NV) & M1 MF1 & M2 MF1 & M3 MF1\\
\midrule
非空 & 126 (126/0) & n/a & n/a & n/a\\
空 & 34 (2/32) & 0.512 & 0.512 & 0.494\\
\bottomrule
\end{tabular}
\end{table}

非空层是单类别，macro-F1 与 balanced accuracy 不可定义；报告该层时只能给出正类 F1/召回，不能填 0。空层仅 34 条，区间很宽，M3 的差值也不显著，因此该分层只作混淆检查，不作因果结论。

\subsection{配对测试集上的检测与联合判别}
表~\ref{tab:pair-detection} 给出三种策略在配对测试集上的分侧检测结果：漏洞侧 180 条、修复侧 180 条、总计 360 条函数记录。漏洞侧的精确率、召回率和 F1 以有漏洞为正类；修复侧的准确率和特异度衡量模型能否接受修复版本；总体准确率把 360 条函数视为独立观测，仅用于描述。三种策略的漏洞侧精确率都约为 0.56--0.59，但召回率只有 0.36--0.39；修复侧准确率 71.7\%--74.4\%。覆盖率均为 100\%。

\begin{table*}[t]
\caption{配对测试集上的分侧检测结果（180 对 / 360 条函数）}
\label{tab:pair-detection}
\centering
\small
\begin{tabular}{@{}lccccccc@{}}
\toprule
策略 & 漏洞 P & 漏洞 R & 漏洞 F1 & 修复 Acc & 修复 Spec & 总体 Acc & Cov.\\
\midrule
M1 零样本直接判定 & 0.586 & 0.361 & 0.447 & 74.4 & 74.4 & 55.3 & 100.0\\
M2 固定少样本 & 0.560 & 0.361 & 0.439 & 71.7 & 71.7 & 53.9 & 100.0\\
M3 上下文增强 & 0.587 & 0.394 & 0.472 & 72.2 & 72.2 & 55.8 & 100.0\\
\bottomrule
\end{tabular}
\end{table*}

表~\ref{tab:pair-joint} 给出配对联合结果。正确区分是漏洞侧判 1、修复侧判 0；双报漏洞表示模型对修复不敏感；双报安全表示模型漏掉漏洞；漏报加误报表示方向反了。联合判别准确率（DA）以 180 对为分母；valid 与 all 版本在本数据上都等于 180/180，因为两侧都有有效输出。DA 的机会水平见表~\ref{tab:ref-paired}：常量基线为 0\%，配对内随机指定一侧为 56.1\%（解析期望 50\%）。三种策略的 DA 只有 8.9\%--12.2\%，显著低于随机机会；主要失败是双报安全（108--113 对）和双报漏洞（44--49 对），正确区分只有 16--22 对。

\begin{table*}[t]
\caption{配对联合判别结果（180 对）}
\label{tab:pair-joint}
\centering
\small
\begin{tabular}{@{}lccccccc@{}}
\toprule
策略 & 正确区分 & 双报漏洞 & 双报安全 & 漏报+误报 & DA(valid) & DA(all) & 有效对/Cov.\\
\midrule
M1 零样本直接判定 & 21 & 44 & 113 & 2 & 0.117 & 0.117 & 180/180 (100.0)\\
M2 固定少样本 & 16 & 49 & 113 & 2 & 0.089 & 0.089 & 180/180 (100.0)\\
M3 上下文增强 & 22 & 49 & 108 & 1 & 0.122 & 0.122 & 180/180 (100.0)\\
\bottomrule
\end{tabular}
\end{table*}

\subsection{修复生成的结构结果}
表~\ref{tab:repair} 汇总修复任务的结构指标。修复请求共 308 条：独立测试集 128 条漏洞版本，配对测试集 180 条漏洞版本。schema 有效率、完整函数率、与输入不同率、完整且不同率以 308 条或 schema 有效输出为分母；精确匹配和空白折叠匹配只对配对测试集的 180 条有参考修复的漏洞版本计算。独立测试集的 128 条没有参考修复，不能计入匹配率的分母。few-shot 的 schema 有效率 99.0\% 最高，direct 84.4\%，context 83.4\%；完整且不同率分别为 80.0\%、62.0\%、82.9\%。精确匹配参考修复只有 1.7\%、5.0\%、2.2\%，空白折叠匹配为 1.7\%、5.0\%、3.9\%。这些指标都是结构代理，不是“修复正确率”；文本匹配低不代表补丁错，匹配高也不代表补丁对。

\begin{table*}[t]
\caption{修复生成的结构代理结果（308 条测试漏洞版本）}
\label{tab:repair}
\centering
\small
\begin{tabular}{@{}lccccccc@{}}
\toprule
策略 & Schema & 完整函数 & 与输入不同 & 完整且不同 & 精确匹配 & 空白折叠 & 参考数\\
\midrule
M1 零样本直接判定 & 84.4 & 99.6 & 80.4 & 80.0 & 1.7 & 1.7 & 180\\
M2 固定少样本 & 99.0 & 99.0 & 62.6 & 62.0 & 5.0 & 5.0 & 180\\
M3 上下文增强 & 83.4 & 99.6 & 83.3 & 82.9 & 2.2 & 3.9 & 180\\
\bottomrule
\end{tabular}
\end{table*}

“完整且不同”是最低限度的可用性检查：模型至少返回了一个完整函数，而不是原样复制或只给 diff。它不能证明漏洞被消除；一个把 \texttt{memcpy} 换成 \texttt{strcpy} 的补丁也可能“完整且不同”，但引入了新的风险。精确匹配参考修复率同样不是正确率：参考修复只是开发者的一种写法，模型可以用不同但语义等价的写法修复；文本匹配低不代表修复错，文本匹配高也不代表修复对。报告在定性分析中会给出“完整且不同但语义可疑”的案例。所有修复结论只能限定为结构可用性；人工复核即使执行，也只报告 rubric 下的通过计数和评分者一致性，不换算成准确率。

\subsection{方法间差异与统计检验}
表~\ref{tab:mcnemar} 记录检测与配对任务的三组两两 McNemar 比较。修复任务没有做 McNemar：结构指标的分母和有效样本随策略不同，且“完整且不同”不是正确性标签；修复质量交给第~\ref{sec:human-review}节的盲评。报告原始 $p$ 与 Holm 校正后的 $p$，并给出 discordant 数。检测任务与配对任务共 6 组比较，在 Holm 校正后都不显著；配对任务中 M1 vs.\ M2 的原始 $p=0.2266$（Holm $0.4531$）相对较小，仍不显著。由于同项目/同对单元相关，McNemar 只作次要证据，主推断使用 10{,}000 次 project/pair cluster bootstrap 差值区间。

\begin{table}[t]
\caption{策略间 McNemar 检验（双侧）}
\label{tab:mcnemar}
\centering
\scriptsize
\begin{tabular}{@{}llccc@{}}
\toprule
任务 & 比较 & A对B错 / A错B对 & 原始 $p$ & Holm $p$\\
\midrule
检测 & M1 vs.\ M2 & 4 / 7 & 0.5488 & 1.0000\\
检测 & M1 vs.\ M3 & 4 / 4 & 1.0000 & 1.0000\\
检测 & M2 vs.\ M3 & 6 / 3 & 0.5078 & 1.0000\\
联合 & M1 vs.\ M2 & 8 / 3 & 0.2266 & 0.4531\\
联合 & M1 vs.\ M3 & 5 / 6 & 1.0000 & 1.0000\\
联合 & M2 vs.\ M3 & 3 / 9 & 0.1460 & 0.4380\\
\bottomrule
\end{tabular}
\end{table}

检测任务的 macro-F1 差值（A−B）为：M1−M2 $-0.0167$（95\% 区间 $-0.0581$--$0.0119$）、M1−M3 $0.0015$（$-0.0328$--$0.0364$）、M2−M3 $0.0182$（$-0.0102$--$0.0657$）；所有区间都跨 0。配对联合准确率的差值为：M1−M2 $0.0278$（$-0.0056$--$0.0667$）、M1−M3 $-0.0056$（$-0.0390$--$0.0278$）、M2−M3 $-0.0333$（$-0.0722$--$0.0056$），同样全部跨 0。因此不能把 few-shot 的数值最高解释为稳定优势。

\subsection{研究问题的回答}
\textbf{RQ1.} 三种策略在独立测试集上的 macro-F1 分别为 0.383、0.399 和 0.381，MCC 分别为 0.041、0.061 和 0.020，特异度分别为 0.750、0.750 和 0.719。few-shot 的 macro-F1 数值最高，但 95\% 区间高度重叠（0.305--0.593 对 0.305--0.552、0.286--0.560），两两差值区间全部跨 0，McNemar Holm 后均不显著。结论只适用于 160 条冻结测试记录和当前基础模型；多数类基线为 80.0\% 准确率、0 特异度、0 MCC，不能把接近 80\% 的准确率当作有效检测。

\textbf{RQ2.} 三种策略的联合判别准确率（有效对）分别为 0.117、0.089 和 0.122，95\% 区间分别为 0.078--0.172、0.056--0.140、0.082--0.178。正确区分、双报漏洞、双报安全、漏报加误报的计数分别为 21/44/113/2、16/49/113/2、22/49/108/1。测试集随机指定一侧的双方均正确率为 56.1\%（101/180，解析期望 50\%），三种策略都远低于这一机会水平；主要失败是双报安全（108--113 对）和双报漏洞（44--49 对）。配对联合指标不能由独立集函数级指标推出，二者排序也不一致。

\textbf{RQ3.} 三种策略的 schema 有效率分别为 84.4\%、99.0\% 和 83.4\%，完整且不同率分别为 80.0\%、62.0\% 和 82.9\%，精确匹配参考修复率分别为 1.7\%、5.0\% 和 2.2\%。这些都是结构指标，不是修复正确率；few-shot 的 schema 最高但“与输入不同”只有 62.6\%，说明格式合规不等于补丁有实质改动。LLM 辅助盲评的 VR=2 为 17/81（21.0\%），direct 7/26（26.9\%）、context 6/25（24.0\%）、few-shot 4/30（13.3\%）；盲评不是人类评审，未编译、未运行 PoC，MIN 的一致性极低，因此只能作为探索性证据。

\textbf{RQ4.} 三种策略的总费用分别为 0.646、1.131 和 0.662 美元（每 100 个评估单元 0.190、0.333、0.195 美元），p95 延迟分别为 19.06、17.87、19.99 秒。few-shot 的费用最高，主要来自更长的少样本输入；allowance 是 OpenCode Go 套餐额度，分别为 0.191、0.183、0.195，不能与 USD 成本相加。成本与效果的联合解读见第~\ref{sec:cost}节。

\textbf{RQ5.} 10 个正式案例覆盖独立集误报/漏报、策略分歧、配对双报漏洞/双报安全/判反、修复 schema 无效、完整但参考不匹配以及 context 高/低风险判定；精炼表见第~\ref{sec:qualitative}节。案例用于展示错误类型，不用于估计总体错误率。

\section{成本}
\label{sec:cost}

\subsection{计费口径与构成}
计费公式见第~\ref{sec:eval}节。成本不仅包括 token：提示长度、输出长度、重试率和延迟都会影响实际可用性。M2 的输入比 M1 多 4 条检测示例或 2 条修复示例，M3 的输入比 M1 多上下文；因此三者的输入 token 不应相同。报告应给出每种策略的平均输入 token、平均输出 token、总调用次数和总费用，而不是只给一个总价。输出 token 的差异还可能反映模型是否遵循“完整函数”的格式：修复任务需要输出整个函数，输出 token 通常远高于检测任务，这部分费用应计入方法成本。

成本按 \path{results.jsonl} 的记录级 \texttt{task} 字段汇总，而不是使用 \path{metrics.json} 的跨任务 usage。每条记录读取 prompt/completion/cached token、\texttt{response.cost.cost\_usd}、\texttt{response.usage.opencode.allowance} 和 \texttt{response.request\_latency\_seconds}。\texttt{allowance} 是 OpenCode Go 套餐额度，不是 USD 成本，必须与 \texttt{cost\_usd} 分述，不能相加；表中的“总 USD”只含 \texttt{cost\_usd}，\texttt{allowance} 在 JSON 和填表脚本注释行中单独给出。只有全部计划记录都存在且每条记录都有数值成本时才给出总成本，否则记 \texttt{null} 或 \texttt{---}；schema-invalid 与 API-error 分别计数；表中的“无效率 (schema/API)”分别给出两者的百分比。“每100 USD”以 340 个评估单元（160 条检测函数 + 180 对）为分母，不是每 100 次请求。

除货币成本外，报告 p50 与 p95 延迟、总墙钟时间和失败重试次数。延迟在共享网络和限流条件下波动较大，因此以分位数描述，并注明测量地点和并发度。若实验在本地或私有部署上运行，则用 GPU 小时和折旧估算成本，不能直接套用 API 单价。

\begin{table*}[t]
\caption{三种策略的调用成本与无效输出（每策略 828 次请求）}
\label{tab:cost}
\centering
\scriptsize
\begin{tabular}{@{}lccccccc@{}}
\toprule
策略 & 检测 in/out & 修复 in/out & 总 USD & 每100 USD & p50 (s) & p95 (s) & 无效率 (schema/API)\\
\midrule
M1 零样本直接判定 & 492{,}158/48{,}053 & 283{,}000/296{,}145 & 0.645585 & 0.1899 & 6.20 & 19.06 & 5.8/0.0\\
M2 固定少样本 & 1{,}301{,}278/34{,}377 & 1{,}172{,}812/289{,}920 & 1.131383 & 0.3328 & 5.80 & 17.87 & 0.4/0.0\\
M3 上下文增强 & 528{,}578/47{,}143 & 305{,}069/296{,}321 & 0.662251 & 0.1948 & 6.25 & 19.99 & 6.2/0.0\\
\bottomrule
\end{tabular}
\end{table*}

三种策略的 OpenCode Go allowance 分别为 0.191327、0.182808、0.195458，均与 USD 成本分开记录，不进入“总 USD”。few-shot 的输入 token 最多（检测 1{,}301{,}278、修复 1{,}172{,}812），总费用 1.131 美元，约为 direct 的 1.75 倍；其总 schema 无效率最低（0.4\%），修复任务为 0.97\%、direct 为 15.6\%、context 为 16.6\%。context 的总无效率 6.2\% 主要来自修复任务（51/308）；direct 的 5.8\% 也主要来自修复（48/308）。

\subsection{归一化成本口径：configured 与 provider-list 分列}
\label{sec:cost-normalized}
\path{analysis/cost_normalization/} 只读取已保存的 \path{results.jsonl} 与 manifest，按每条 usage 的未缓存输入、缓存命中、输出 token 重算，不修改原始 \texttt{cost\_usd}，也不调用 API。表~\ref{tab:cost-basis} 把四个口径分列：\textbf{configured} 是原始 run 记录的成本；\textbf{normalized} 是按 provider 非高峰 list 价重算的值（DeepSeek 0.15/0.60/0.003，MiMo 0.14/0.28/0.0028 USD per 1M）；\textbf{provider list} 是 \texttt{usage.opencode.list\_cost} 的逐请求求和；\textbf{allowance} 是 provider 侧的套餐额度扣减。allowance 既没有从 list cost 扣除，也没有加到 configured 上，因此不与前三列相加；它不是 USD 账单。

DeepSeek 的 configured 0.30/1.20/0.006 是高峰价，而本批 9{,}584 条有 usage 的请求全部落在非高峰窗口（peak rows = 0），provider 快照使用 0.15/0.60/0.003，所以 cache 口径一致的 v1 run 上 configured 恰为 normalized 的 2 倍；v2 已直接配置为非高峰价，两列相同。71 个完整 run 的 configured 合计 3.886970 美元、normalized 2.146513 美元；17 个 headline run 为 0.558073/0.365704。另有 2 个未完成 run（52 次请求）configured 与 normalized 均为 0.008197，单列、不并入 complete-only 合计；6 个无 usage 的失败/空跑不补零。hybrid v1（10 run、1{,}400 请求）为 0.173865/0.114598，v2 为 0.082335/0.082335；两者请求数与 prompt 几乎相同，USD 差异来自价格口径而非工作量。MiMo 无峰谷价，configured 与 normalized 相同；Xiaomi 端点的 \path{juliet_llm} 不返回 provider list/allowance，按 n/a 处理，不补零。

修复闭环的 all-in 成本修正为 0.082550 美元（configured；normalized 0.052549），不是旧 README 的 0.077516 美元；差额 0.005034 来自解析器开发 run \texttt{parser\_dev/mimo\_test70\_v1} 的真实已发生成本（44/70 单元，不进入 headline）。两个 headline run 仍为 0.032831 美元（normalized 0.019336），不变。DeepSeek 每个 run 的 allowance 约为其 list cost 的 0.600 倍，MiMo 为 1.000 倍；它是额度而非账单，故只在表中单列。


% Auto-generated by report/scripts/build_extended_tables.py. Do not edit.

\begin{table*}[t]
\caption{成本口径：configured 与 normalized provider-list 分列（USD；allowance 为额度扣减，不相加）}
\label{tab:cost-basis}
\centering
\scriptsize
\setlength{\tabcolsep}{3.6pt}
\begin{tabular}{@{}lrrrrrr@{}}
\toprule
组 & run & 调用 & configured & normalized & provider list & allowance\\
\midrule
全部完整 run & 71 & 9,532 & 3.886970 & 2.146513 & 2.123231 & 1.624966 \\
headline 17 run & 17 & 2,362 & 0.558073 & 0.365704 & 0.361279 & 0.284329 \\
hybrid v1（审计） & 10 & 1,400 & 0.173865 & 0.114598 & 0.114599 & 0.090891 \\
hybrid v2（主结果） & 10 & 1,400 & 0.082335 & 0.082335 & 0.082335 & 0.059878 \\
juliet\_llm 修复后两次 & 2 & 280 & 0.006262 & 0.006262 & --- & --- \\
修复 headline & 2 & 142 & 0.032831 & 0.019336 & 0.019336 & 0.013938 \\
修复 all-in & 7 & 355 & 0.082550 & 0.052549 & 0.052549 & 0.040550 \\
\bottomrule
\end{tabular}
\end{table*}



\subsection{成本—效果权衡}
few-shot 以最高成本换取了数值最高的独立集 macro-F1（0.399 对 0.383/0.381），但差值区间跨 0，不能证明这笔额外费用带来可靠收益。direct 最便宜，修复 schema 无效率却最高；context 的成本和延迟接近 direct，配对联合准确率数值最高（0.122），同样不显著。若目标是低成本筛查，direct 的单位成本最低；若目标是减少修复格式失败，few-shot 的 99.0\% schema 有效率有优势，但它的“与输入不同”率只有 62.6\%，且盲评 VR=2 最低（4/30），不能据此推荐自动修复。每正确检测/联合判别的费用未单独列表；在联合准确率仅 8.9\%--12.2\% 的情况下，该指标会被极低分母放大，不适合作为主要排名依据。

\section{定性分析与误差案例}
\label{sec:qualitative}

\subsection{案例抽样规则}
定性分析不使用随机挑选的“好看”案例，而是由 \path{qualitative/cases.json} 按预定义的 10 个类别各抽取 1 个案例：独立集误报、独立集漏报、策略分歧、配对双报漏洞、配对双报安全、配对判反、修复 schema 无效、修复完整但参考不匹配、context 高风险判定、context 低风险判定。选择规则版本为 \texttt{qualitative-v1}，所有输入为六次完整运行；脚本不调用 API。案例编号、原始索引、配对 ID 和代码摘录写入 JSON/Markdown，正文只保留精炼事实，避免泄露无关项目信息。案例解释必须引用函数中的具体语句或数据流，不能只复述模型理由；这些案例用于展示错误类型，不用于估计总体错误率。

\subsection{误差分类框架}
本文按以下维度归纳误差，而不是简单罗列“模型答错”：
\begin{itemize}
    \item \textbf{证据缺失}：函数确实存在缺陷，但提示中没有出现触发条件或调用者信息，模型无法从可见代码推断；
    \item \textbf{证据误读}：模型引用了不存在的变量、错误的边界条件或与代码相反的返回值语义；
    \item \textbf{表面启发式}：模型因函数名、危险 API、注释或代码长度而判定，未检查实际数据流；
    \item \textbf{上下文误导}：M3 因上下文不完整或标签相关的“有无上下文”格式而改变判断；
    \item \textbf{少样本模仿}：M2 机械模仿示例的 CWE、理由或代码模式，忽略当前函数的数据流；
    \item \textbf{配对失败}：双报漏洞、双报安全或方向反转，说明模型没有利用修复差异；
    \item \textbf{修复症状化}：只改表面语句、加长缓冲区或改注释，未处理根因；
    \item \textbf{修复语义改变}：补丁改变了返回值、错误处理或 API 约定；
    \item \textbf{修复新风险}：引入新的不安全调用、缺少检查或更宽的攻击面；
    \item \textbf{格式与解析}：输出不符合 JSON 模式、缺少完整函数、被截断或自相矛盾。
\end{itemize}
前五类主要反映检测能力，第六类反映配对判别，第七至第九类反映修复质量，第十类反映工程实现。报告中应分别统计各类的比例，避免把解析失败与模型判断错误混为一谈。

\subsection{案例表}
表~\ref{tab:cases} 为 10 个正式案例的精炼表，完整代码摘录、逐策略输出和选择理由见 \path{qualitative/cases.md} 与 \path{qualitative/cases.json}。案例来自预定义分层，覆盖独立集误报/漏报、策略分歧、配对三类失败、修复格式/文本失败和 context 高/低风险；它们用于展示错误类型，不用于估计总体错误率。

\begin{table*}[t]
\caption{10 个误差案例的精炼表（完整证据见 \texttt{qualitative/cases.md}）}
\label{tab:cases}
\centering
\scriptsize
\begin{tabular}{@{}l>{\raggedright\arraybackslash}p{2.6cm}>{\raggedright\arraybackslash}p{2.9cm}>{\raggedright\arraybackslash}p{4.3cm}>{\raggedright\arraybackslash}p{5.2cm}@{}}
\toprule
编号 & 类别 & 案例 & 关键事实 & 解释\\
\midrule
C-01 & 独立集误报 & idx 8632 & 3/3 判漏洞，平均置信度 0.917；标签非漏洞 & 模型把 eXIf 分支的未初始化读/边界推理当成漏洞；函数级标签与模型证据不一致，需人工复核 \\
C-02 & 独立集漏报 & idx 13796 & 3/3 判非漏洞，平均置信度 0.950；标签漏洞（CWE-125/787） & 触发路径未在片段中出现，模型被局部边界检查误导，属于漏报 \\
C-03 & 策略分歧 & idx 2548 & direct 非漏洞 0.600；few-shot 漏洞 0.900；context 漏洞 0.750 & 示例/上下文改变判断，置信度跨度 0.300；同一函数上策略不稳定 \\
C-04 & 配对双报漏洞 & 19923:19924 & 3/3 给出 both\_flagged；两侧平均置信度 0.908 & 模型对修复不敏感，把修复侧也判为漏洞 \\
C-05 & 配对双报安全 & 25144:25145 & 3/3 给出 both\_clean；两侧平均置信度 0.947 & 漏洞侧和修复侧都判安全，属于漏报 \\
C-06 & 配对判反 & 15922:15928 & 1/3 给出 reversed；direct/context 为 both\_flagged & 方向反转：漏洞侧判安全、修复侧判漏洞 \\
C-07 & 修复 schema 无效 & 10480:10481 & direct 输出缺少 \texttt{explanation}，校验失败 & 格式/字段失败，不是补丁质量；该错误在无效修复记录中出现 93 次 \\
C-08 & 修复完整但参考不匹配 & 25161:25163 & context 输出完整且不同；行相似度 0.605，差异 47 行 & 完整且不同不等于正确；文本不匹配也不能证明补丁错，需语义/人工审查 \\
C-09 & context 高风险判定 & 679:680 漏洞侧 & context 非空，判漏洞 0.900，标签漏洞 & 预测正确，但独立集 context 非空层与标签高度相关，不能当作上下文因果证据 \\
C-10 & context 低风险判定 & 16915:16937 修复侧 & context 非空，判非漏洞 0.900，标签非漏洞 & 预测正确，但 context 分层是混淆检查，不是因果证据 \\
\bottomrule
\end{tabular}
\end{table*}

\subsection{分析步骤}
案例由定性脚本从真实运行按上述类别导出，报告作者依据代码摘录、标签和逐策略输出整理解释；这不是双人独立人工审计，也没有案例层面的评分者一致性。模型给出的自然语言理由只作为“模型声称的证据”，必须与代码对照；不能把模型理由直接当作案例解释。对于涉及上游项目历史的问题，只使用本地数据中可见的字段，不联网检索修复提交，以免引入评估时不可见的信息。

\subsection{修复结果的 LLM 辅助盲评}
\label{sec:human-review}
修复盲评已完成，产物在 \path{repair_eval/analysis/repair_blind_v1/final/}（包 ID \texttt{PKG-3DB061ED2788}），rubric 版本 \texttt{repair-eval-rubric-2026-11-11.1}。盲评包含 30 个案例、81 个候选；候选只包含 schema 有效的修复输出，因此 schema 无效记录不在盲评框内。评审为两名 LLM 辅助评分者加第三 LLM 裁决，\textbf{不是人类评审}；汇总文件的裁决者身份未记录。评分者看不到策略名、模型名、pair/sample ID、项目和路径；候选 ID 由私有 salt 和 HMAC-SHA256 生成。所有 81 个候选的 compile\_status 都是 \texttt{NOT\_RUN}：没有编译、测试或 PoC oracle。

评分使用七维 rubric，每维 0/1/2：漏洞消除（VR）、行为保持（SP）、最小性（MIN）、完整性（COM）、代码质量（CQ）、解释忠实（EF）、无新风险（NS）。COM 是文本完整性，不是编译证明；CQ 是可读性，不是漏洞是否消除。主通过条件是 VR=2，严格通过条件是 VR=2 且 SP$\ge$1、NS=2、EF$\ge$1、COM$\ge$1。最终结果：VR=2 为 17/81（21.0\%），严格通过同为 17/81；按策略 direct 7/26（26.9\%）、context 6/25（24.0\%）、few-shot 4/30（13.3\%）。七维均值为 1.547（仅供参考）。

\begin{table*}[t]
\caption{修复盲评：七维分数与评分者一致性（30 案例 / 81 候选；LLM 辅助评审，不是人类）}
\label{tab:blind}
\centering
\small
\begin{tabular}{@{}lcccc@{}}
\toprule
维度 & 均值 & 2 分率 & 二次加权 $\kappa$ & 完全一致率\\
\midrule
VR 漏洞消除 & 0.481 & 0.210 & 0.902 & 0.901\\
SP 行为保持 & 1.765 & 0.778 & 0.409 & 0.802\\
MIN 最小性 & 1.593 & 0.593 & 0.015 & 0.580\\
COM 完整性 & 2.000 & 1.000 & 0.000 & 0.988\\
CQ 代码质量 & 1.963 & 0.963 & 0.654 & 0.975\\
EF 解释忠实 & 1.148 & 0.259 & 0.658 & 0.679\\
NS 无新风险 & 1.877 & 0.914 & 0.621 & 0.926\\
\bottomrule
\end{tabular}
\end{table*}

表~\ref{tab:blind} 汇总七维分数和评分者一致性。MIN 的二次加权 $\kappa$ 只有 0.015，说明两名评分者对“最小性”几乎没有一致意见~\cite{landis1977}；COM 的 $\kappa=0.000$ 是因为 81 个候选全部为 2 分，分数恒定，$\kappa$ 不可解释。其他维度的 $\kappa$ 也多为中低水平，因此不能把 17/81 当作稳定的人类判断。事后结构检查（不进入 rubric）显示：三个策略的完整函数率都是 100\%；有改动率 context 96.0\%、direct 88.5\%、few-shot 53.3\%；参考文本精确匹配率均为 0\%；平均改动行数分别为 8.0、7.3、1.9。few-shot 的“有改动率”最低，和它较低的 VR=2 比例一致。

这些结果只能支持“在该七维 rubric 下，LLM 辅助评审认为 17/81 个候选处理了标注 CWE 的根因”。它们不证明补丁可编译、通过测试、保持行为或消除漏洞；候选集只含 schema 有效输出，每策略候选数不同（direct 26、context 25、few-shot 30），不能把候选比例直接当作策略总体质量。MIN 的低一致性和 LLM 评审的性质使结论保持探索性。

\section{讨论}
\label{sec:discussion}

\subsection{三种评估视图的互补性}
独立集回答“给定一条函数，模型能否识别漏洞”，配对集回答“同一个函数被修复后，模型能否同时判对两个版本”，修复任务回答“模型能否生成一个结构上可用的补丁”。三者都不是同一个能力的简单重复：独立集受类别比例和样本构成影响；配对集受配对内部相关性影响；修复任务受输出长度、格式和参考写法的限制。一个在独立集上精确率很高、在配对集上双报安全比例也很高的模型，可能只是倾向于判“无漏洞”；一个在配对集上 DA 很高、但修复任务原样返回的模型，可能只学会了“保守判定”，没有学会“生成补丁”。报告中把三组指标并列，而不是合并成一个总分，是为了让读者看到这种权衡。

配对集还有一个容易被忽略的性质：同一对的两个版本共享项目、文件、函数名和大部分代码，只有修复相关的变化不同。因此 DA 高并不必然意味着模型理解了修复；它可能只是对代码的细微变化不敏感，或对两个版本都给出同一标签。表~\ref{tab:pair-joint}同时报告四类联合结果，正是为了区分“真的区分”与“碰巧答对”。在解释结果时，应优先检查正确区分和双报漏洞的比例，而不是只看版本级准确率。

\subsection{策略差异的来源}
三种策略的差异可以分解为三个通道：示例、上下文和提示长度。M2 只增加示例，因此若 M2 优于 M1，差异更可能来自任务格式的示范和对标签边界的校准；M3 增加上下文，若 M3 优于 M1，差异可能来自跨函数信息，也可能来自“有没有上下文”的格式差异。检测测试集的 context 非空层 126/126 为漏洞样本，空层 34 条中 2 条漏洞、32 条非漏洞；非空层是单类别，M3 的总体 macro-F1（0.381）与 M1（0.383）基本持平，空层 M3 为 0.494、M1 为 0.512，差值 $-0.0179$（区间 $-0.1071$--0.0000）。因此没有证据表明 M3 的上下文带来了可靠提升，M3 的结果不能解释为“模型理解了项目上下文”。

另一个通道是提示长度与输出格式。三种策略使用相同的英文任务说明和 JSON 模式，但 M2、M3 的附加内容会改变模型对任务的表征。更长的提示可能让模型更谨慎，或更倾向于输出“有漏洞”。为检查这一点，结果部分应报告预测正类率（预测为有漏洞的比例），并检查三种策略的正类率是否接近；若 M3 的正类率显著高于 M1，那么 F1 的变化可能只是判定阈值的移动，而不是判别能力的提升。定性分析应抽取策略分歧样本，检查模型是否引用了示例中的具体模式，或是否把上下文中的外部函数名当作危险信号。

\subsection{修复结构指标的解释边界}
修复任务的结构指标只能回答“输出是否像一个补丁”，不能回答“补丁是否正确”。完整函数率低说明模型没有遵循格式；与输入不同率低说明模型可能原样返回；完整且不同率高说明输出至少可用，但仍可能只修了症状、改变了语义或引入新风险。精确匹配参考修复率衡量的是文本相似度，不是语义等价：模型可以用不同写法正确修复，也可以用相同写法引入新问题。第~\ref{sec:human-review}节的 LLM 辅助盲评已完成，但评审不是人类，81 个候选的 compile\_status 全部为 NOT\_RUN，MIN 的评分者一致性仅 0.015；因此 17/81 的 VR=2 只能写成“在该 rubric 下被评为根因已处理”，不能换算成修复正确率或漏洞消除率。

在工程实践中，修复生成的门槛应高于检测：检测的假阳性可以由人工复核消化，修复的错误可能直接进入代码库。因此，即使某策略的完整且不同率较高，若人工复核发现语义改变或新风险，也不应推荐自动合并。本文不预设哪种策略适合哪种场景；这些判断应在结果表填写后，结合具体项目的审查流程和风险预算给出。

\subsection{与训练型检测器的可比性}
本文的三种策略都在推理时通过提示调用模型，不更新参数，也不使用项目特定的训练数据。与 Devign、LineVul 等训练型检测器~\cite{zhou2019devign,fu2022linevul}相比，提示方法的优势是部署快、无需标注训练集、可通过修改提示适配新任务；劣势是无法利用项目分布，单次调用成本随样本量线性增长，且行为受服务端模型版本影响。本文不把当前指标与这些论文的报告值直接比较，因为数据划分、标签定义、评估单位和是否微调都不同。若后续需要与训练型基线比较，应在同一划分上重新训练或调用其公开模型，并报告相同的指标族。

\subsection{有效性威胁}
\textbf{构造效度。} 标签来自 SecVulEval，本文不重新审计代码。无漏洞标签表示数据源未标注该函数的漏洞，不证明函数不存在其他缺陷；修复版本也可能仍含其他问题。模型输入只给函数和受限上下文，真实审查中可用的调用链、构建配置和威胁模型可能更多，因此本文结果不能代表信息更充分时的性能上限。检测测试集的 context 非空层 126/126 为漏洞，空层 34 条中只有 2 条漏洞，关联很强且非空层为单类别，这是 M3 最主要的构造威胁；源级 32{,}767 字符截断也使超长函数的结果只能解释为截断后的下界。此外，\texttt{fixed\_func\_idx} 的存在性与标签相关（测试集正例 115/128、负例 0/32；源子集 143/160 与 0/40），CWE/CVE 在 32/32 条测试负例（源子集 40/40）中也非空；协议虽然排除这些字段，但字段分布本身说明数据构造带有强标签信号，任何未按白名单过滤的提示都会虚高。CWE-401 在源子集只有 4 条、测试集 3 条，且严格 C++ 单标签为 0，逐 CWE 结论同样受构造效度限制。

\textbf{内部效度。} 三种策略的提示长度不同，模型可能对长度、示例顺序或上下文格式敏感。温度设为 0 只能减少采样随机性，不能保证服务端完全确定。解析器和重试策略若在策略间不一致，会引入工程偏差；本文要求共用同一实现，并在日志中记录解析状态。修复任务没有编译/测试 oracle，结构指标无法区分“正确修复”和“看似完整的错误修复”；LLM 辅助盲评已完成，但它不是人类评审，也没有编译或 PoC，MIN 的评分者一致性极低，因此修复结论仍是探索性的。开发示例来自冻结开发集，已通过运行器的测试集隔离检查，但示例仍会影响模型的输出风格；few-shot 的结果不能等同于“在更大开发集上调参”的结果。表~\ref{tab:ref-independent} 和表~\ref{tab:ref-paired} 的规则基线不是模型，不能用于估计模型方差或替代缺失的模型结果。

\textbf{外部效度。} 源子集共 600 条记录、项目并集 134 个，冻结测试集为 520 条函数记录（160 检测 + 360 配对函数），规模有限；检测集采用 80/20 的人为比例，配对集标签平衡，两者都不代表真实代码库的漏洞 prevalence。项目层面在任务内已隔离，但跨任务仍有 2 个检测开发项目出现在配对测试、5 个配对开发项目出现在检测测试；源子集共享 26 个项目，不能据此声称完全跨项目泛化。模型可能在预训练中见过公开 CVE 及其修复，本文无法完全排除污染；因此结果应理解为“在当前模型和当前数据上的表现”，而不是对未知项目的保证。CWE-401 在测试集只有 3 条、只来自 2 个项目（linux 2 条、ImageMagick 1 条），且 Tier A 的 SecVulEval 数据中 0 条 Juliet 记录；即使模型在这几条上表现良好，也不能外推为 CWE-401 专项、C++ 内存泄漏或 Juliet 基准上的能力。Tier B 的 Juliet 结果只在第~\ref{sec:ext-juliet}节的分层内解释，不能反过来外推到 Tier A 或真实项目。

\textbf{结论效度。} 主要指标是固定阈值下的分类指标和结构指标，未做阈值调优；若模型输出置信度，改变阈值会改变精确率与召回率的权衡。本文不以单一指标排名，而是报告指标族、置信区间和配对检验。样本量较小时，置信区间可能很宽，方法间的小差异不应被解释为稳定优势。规则基线只提供解释锚点（多数类、常量、配对随机和泄漏哨兵），不是模型对照，也不支持“某策略优于基线”的统计结论。


\section{扩展评估 I：统一分层基准与双模型比较}
\label{sec:ext-tier-a}

旧报告的六次运行只覆盖 \texttt{deepseek-v4.1-flash} 的三种提示策略。扩展部分新增了 MiMo-V2.6-Flash 的六个全量 run、180 对联合提示、160 条证据提示、Juliet CWE-401 的静态/动态/LLM/混合运行，以及 70 个 bad 单元的修复闭环。这些结果的标签来源、评估单元、类别先验、分母和成本口径都不同：Tier A 是真实项目函数片段，Tier B 是合成单 CWE 用例；Tier A 的负例是“源数据未标注漏洞”，Tier B 的 good 是模板生成的修复版本。把它们池化成一个 leaderboard 会产生没有解释的“平均 F1”。因此本节先定义分层契约，再报告 Tier A 的双模型结果。

\subsection{分层契约}
表~\ref{tab:bench-contract} 给出两个层级的边界。Tier A 的主终点是检测 macro-F1 与配对 DA；Tier B 的主终点是泄漏 macro-F1 与 MCC，并强制同时报告 coverage、strict 和失败分布。两层共享以下字段与规则：
\begin{itemize}
    \item 每个单元一行，保留 \texttt{status}（ok / invalid / compile\_error / inconclusive / timeout / tool\_error），不删除失败单元，也不把失败改写成 $y_{\mathrm{pred}}=0$；
    \item 指标同时给 valid-only 与 strict all-record 两个口径；coverage 的分母是计划单元数；
    \item 比例用 Wilson 区间~\cite{wilson1927}；macro-F1、MCC、balanced accuracy 用 10{,}000 次 cluster bootstrap~\cite{efron1979}，检测按项目、配对按 \texttt{pair\_id}、Juliet 按 family 聚类；
    \item 配对正确性用 McNemar 精确检验~\cite{mcnemar1947}，族内用 Holm 校正~\cite{holm1979}；McNemar 只作次要证据，主推断用 cluster bootstrap 差值区间；
    \item 成本分 API USD 与本地计算两类；本地静态/动态工具的 \texttt{usd=null}，不是 0；API 成本按记录级 usage 与价格快照计算，\texttt{allowance} 单独报告。
\end{itemize}
这些规则不保证结论正确，但保证“失败不被静默丢掉、分母不被事后改变、成本不被混算”。表~\ref{tab:bench-contract} 的最后一列明确写出每层不能声称什么，这是与单一排行榜最重要的区别。


% Auto-generated by report/scripts/build_extended_tables.py. Do not edit.

\begin{table*}[t]
\caption{统一分层基准的两个层级与不可合并的原因}
\label{tab:bench-contract}
\centering
\scriptsize
\setlength{\tabcolsep}{3.5pt}
\begin{tabular}{@{}p{1.3cm}p{4.0cm}p{4.6cm}p{2.7cm}p{4.2cm}@{}}
\toprule
层级 & 数据与单元 & 方法 & 主终点 & 明确不外推\\
\midrule
Tier A & SecVulEval 冻结测试集；160 条函数、180 对版本；C/C++ 真实项目抽取片段 & DeepSeek/MiMo $\times$ direct / few-shot / context；配对联合提示；证据提示 & 检测 macro-F1；配对 DA & 不能外推到 Juliet 的合成单 CWE；不报告 Tier A 静态 F1\\
Tier B & Juliet CWE-401 v1.3 冻结子集；200 单元（dev 60 / test 140），100 组、42 家族 & Clang SA、Cppcheck、ASan/LSan、Valgrind、LLM、静态/动态证据混合 & 泄漏 macro-F1 + MCC，强制 coverage/strict & 不能外推到真实项目；动态 clean 不是 TN；单 CWE、模板化\\
\bottomrule
\end{tabular}
\end{table*}


\begin{table*}[t]
\caption{Tier A 检测：DeepSeek 与 MiMo 的三策略结果（160 条：128 正 / 32 负）}
\label{tab:ea-det}
\centering
\scriptsize
\setlength{\tabcolsep}{3.5pt}
\begin{tabular}{@{}llrrrrrrrr@{}}
\toprule
模型 & 策略 & 有效 & P & R & Macro-F1 [95\% CI] & MCC & Strict & 任务 USD & p95 (s)\\
\midrule
DeepSeek & direct & 160/160 & 0.826 & 0.297 & 0.383 [0.305, 0.552] & 0.041 & 0.388 & 0.066988 & 10.83 \\
DeepSeek & few-shot & 160/160 & 0.837 & 0.320 & 0.399 [0.305, 0.593] & 0.061 & 0.406 & 0.136741 & 9.69 \\
DeepSeek & context & 160/160 & 0.812 & 0.305 & 0.381 [0.286, 0.560] & 0.020 & 0.388 & 0.070308 & 10.71 \\
MiMo & direct & 150/160 & 0.806 & 0.242 & 0.346 [0.264, 0.528] & 0.008 & 0.325 & 0.024397 & 12.00 \\
MiMo & few-shot & 160/160 & 0.773 & 0.266 & 0.346 [0.262, 0.512] & -0.042 & 0.350 & 0.034011 & 11.03 \\
MiMo & context & 153/160 & 0.861 & 0.252 & 0.365 [0.278, 0.571] & 0.080 & 0.350 & 0.025568 & 8.38 \\
\bottomrule
\end{tabular}
\end{table*}


\begin{table*}[t]
\caption{Tier A 配对联合判别：同一 180 对上的分侧调用结果}
\label{tab:ea-pair}
\centering
\scriptsize
\setlength{\tabcolsep}{4pt}
\begin{tabular}{@{}llrrrrrrrr@{}}
\toprule
模型 & 策略 & 正确/180 & DA(valid) [95\% CI] & DA(all) & 双报漏洞 & 双报安全 & 判反 & 有效/180 & 任务 USD\\
\midrule
DeepSeek & direct & 21 & 11.7 [7.8, 17.2] & 11.7 & 44 & 113 & 2 & 180/180 & 0.138323 \\
DeepSeek & few-shot & 16 & 8.9 [5.5, 14.0] & 8.9 & 49 & 113 & 2 & 180/180 & 0.294895 \\
DeepSeek & context & 22 & 12.2 [8.2, 17.8] & 12.2 & 49 & 108 & 1 & 180/180 & 0.144837 \\
MiMo & direct & 14 & 9.1 [5.5, 14.7] & 7.8 & 27 & 108 & 5 & 154/180 & 0.048290 \\
MiMo & few-shot & 17 & 9.4 [6.0, 14.6] & 9.4 & 40 & 117 & 6 & 180/180 & 0.067880 \\
MiMo & context & 19 & 11.4 [7.5, 17.2] & 10.6 & 38 & 106 & 3 & 166/180 & 0.049967 \\
\bottomrule
\end{tabular}
\end{table*}


\begin{table*}[t]
\caption{Tier A 同策略、同任务的模型比较（DeepSeek $-$ MiMo；McNemar + Holm）}
\label{tab:ea-modelcmp}
\centering
\scriptsize
\setlength{\tabcolsep}{4pt}
\begin{tabular}{@{}llrrrrrr@{}}
\toprule
任务 & 策略 & $\Delta$ & 95\% CI & $n$ & $b/c$ & 原始 $p$ & Holm $p$\\
\midrule
检测 & direct & 0.031 & [-0.018, 0.077] & 150 & 11/6 & 0.3323 & 0.6646 \\
检测 & few-shot & 0.053 & [0.010, 0.129] & 160 & 15/6 & 0.0784 & 0.2351 \\
检测 & context & 0.009 & [-0.070, 0.058] & 153 & 10/8 & 0.8145 & 0.8145 \\
配对 & direct & 0.032 & [-0.013, 0.078] & 154 & 9/4 & 0.2668 & 0.8005 \\
配对 & few-shot & -0.006 & [-0.039, 0.028] & 180 & 4/5 & 1.0000 & 1.0000 \\
配对 & context & 0.000 & [-0.042, 0.042] & 166 & 6/6 & 1.0000 & 1.0000 \\
\bottomrule
\end{tabular}
\end{table*}


\begin{table*}[t]
\caption{Tier A 模型汇总（configured 口径；描述性；不跨策略池化单元）}
\label{tab:ea-cost}
\centering
\scriptsize
\setlength{\tabcolsep}{3.5pt}
\begin{tabular}{@{}lrrrrrrr@{}}
\toprule
模型 & 检测 macro-F1 均值（范围） & 配对 DA 均值（范围） & 平均覆盖 & 六 run USD & USD/100 请求 & p50/p95 (s) & 请求\\
\midrule
DeepSeek & 0.388 (0.381--0.399) & 0.109 (0.089--0.122) & 1.000 & 2.439219 & 0.098197 & 5.50/10.55 & 2484 \\
MiMo & 0.352 (0.346--0.365) & 0.100 (0.091--0.114) & 0.945 & 0.623154 & 0.025087 & 4.41/10.13 & 2484 \\
\bottomrule
\end{tabular}
\end{table*}



\subsection{Tier A 检测：DeepSeek 与 MiMo}
表~\ref{tab:ea-det} 汇总 6 个检测 run。DeepSeek 的 macro-F1 为 0.383（direct）、0.399（few-shot）、0.381（context）；MiMo 为 0.346、0.346、0.365。DeepSeek 的三个区间分别覆盖 0.305--0.552、0.305--0.593、0.286--0.560；MiMo 为 0.264--0.528、0.262--0.512、0.278--0.571。所有区间高度重叠。MiMo 的 coverage 不是 1.000：direct 有 10 条 schema-invalid（150/160），context 有 7 条（153/160），few-shot 为 160/160；strict accuracy 因此降到 0.325--0.350，而 DeepSeek 为 0.388--0.406。MiMo few-shot 的 MCC 为 $-0.042$，说明其预测与标签在去除类别先验后几乎没有正相关。

检测结果不能脱离类别比例解读。测试集 128 正/32 负，多数类基线为 80.0\% 准确率但特异度 0、MCC 0；DeepSeek/MiMo 的准确率只有 38\%--41\%，看起来“远低于基线”，但其 macro-F1 与 MCC 略高于 0，说明它们至少比常量猜测多识别了一部分正例。反过来，不能因为某个策略的 macro-F1 是 0.399 就宣称它“更好”：区间宽度约 0.25--0.29，0.05 的差距完全落在抽样不确定性内。

\subsection{Tier A 配对：分侧判定的困难}
表~\ref{tab:ea-pair} 给出同一 180 对上的分侧调用结果。DeepSeek 的 valid DA 为 11.7\%、8.9\%、12.2\%，MiMo 为 9.1\%、9.4\%、11.4\%；all-pairs DA 在 MiMo direct/context 上进一步降到 7.8\%/10.6\%，因为至少一侧 invalid 的 pair 记为失败。测试集“随机指定一侧”的双方均正确率为 56.1\%（101/180，解析期望 50\%），所有策略都远低于这一机会水平。

失败结构比 DA 更有信息量。DeepSeek 的 both-safe 为 113/113/108 对，both-vulnerable 为 44/49/49 对，swapped 只有 2/2/1 对；MiMo 的 both-safe 为 108/117/106，both-vulnerable 为 27/40/38，swapped 为 5/6/3。也就是说，模型的主要错误不是“把漏洞侧和修复侧判反”，而是“两侧都判安全”或“两侧都判漏洞”：前者说明模型没有识别出漏洞，后者说明模型对修复不敏感。DA 低并不是因为配对任务“太难随机”，而是因为模型在单侧判定上已经偏向保守/宽松，联合结果被这种偏向主导。

\subsection{模型比较：没有 Holm 显著的差异}
表~\ref{tab:ea-modelcmp} 给出 6 个同策略、同任务的 DeepSeek $-$ MiMo 比较。检测 direct/context 的 $\Delta$ macro-F1 为 0.031/0.009，区间跨 0；few-shot 为 0.053，区间 $[0.010,0.129]$ 排除 0，但 McNemar 原始 $p=0.0784$、Holm $p=0.2351$，不显著。配对 direct/few-shot/context 的 $\Delta$ DA 为 0.032/$-0.006$/0.000，区间全部跨 0。6 个比较经 Holm 后没有一个显著。few-shot 检测出现“bootstrap 区间排除 0、McNemar 不显著”的分歧，原因之一是两者检验的对象不同：bootstrap 对 macro-F1 的配对差值重采样，McNemar 只对逐单元正确性做 discordant 计数；在 160 条、15/6 个 discordant 的情况下，后者功效很低。报告同时给出两者，并明确不把任一单独结果当作排名依据。

模型内策略比较（12 组）也没有 Holm 显著差异：DeepSeek 检测的 direct vs.\ few-shot 为 $-0.017$（$-0.058$--0.012），MiMo 检测的 direct vs.\ few-shot 为 0.008（$-0.016$--0.048）；配对侧同样全部跨 0。MiMo 复制了旧报告的核心结论：在固定模型和单次运行下，三种提示策略没有可检测的稳定优劣。

\subsection{成本与延迟}
表~\ref{tab:ea-cost} 汇总每个模型的六个 run。DeepSeek 2{,}484 次请求花费 2.439219 美元，MiMo 为 0.623154 美元；每 100 次请求分别为 0.098197 与 0.025087 美元；p50/p95 延迟为 5.50/10.55 s 与 4.41/10.13 s。这里的两列 USD 都是 run 记录的 configured 口径：DeepSeek 按 0.30/1.20/0.006 USD per 1M 计价，而本批请求全部落在非高峰窗口，按 provider list 价 0.15/0.60/0.003 归一化后，direct/context 约为 configured 的 0.49--0.50 倍，few-shot 因 legacy 记录不扣缓存而为 0.26--0.29 倍；MiMo 使用 0.14/0.28/0.0028，configured 与 normalized 相同。两套口径不得混用，逐组对照见第~\ref{sec:cost}节表~\ref{tab:cost-basis}。MiMo 更便宜、略快，但其平均 coverage 为 0.945，DeepSeek 为 1.000；按“每有效预测”而不是“每请求”计算，差距会缩小。且 MiMo 需要 \texttt{reasoning\_effort=none} + \texttt{max\_completion\_tokens} 才与 DeepSeek 的非思考模式语义对齐。因此表~\ref{tab:ea-cost} 是“模型 + 传输/解码配置”的比较，不是纯模型能力比较。

\section{扩展评估 II：联合提示与结构化证据}
\label{sec:ext-prompt}

上一节的配对任务使用两次独立判定：模型分别看漏洞版本和修复版本，然后由评测端组合。这样的协议允许模型输出“两侧都安全”或“两侧都漏洞”，也允许它完全忽略另一版本。一个自然的改进是把两个版本放进同一个请求，强制模型二选一。本节报告这个联合提示的真实运行，并随后报告另一个方向的提示：不是给模型看另一版本，而是要求它给出证据行、机制和反事实修复。

\subsection{Pair-joint 协议}
每对只发一次请求。system 说明两段代码是同一 C/C++ 函数的两个修订，恰好一个包含安全漏洞；A/B 顺序由 \texttt{HMAC-SHA256("secvul-pair-order-v1", pair\_id)} 决定，前 90 对漏洞版本在 A、后 90 对在 B，两个模型使用同一映射。用户消息把两个版本按 \texttt{行号 | 源码} 编号，要求返回 \texttt{vulnerable\_version}（A/B）、\texttt{confidence}、\texttt{evidence} 和 \texttt{rationale}；\texttt{evidence} 至少一项，且每项必须引用所选版本的行号和对应源码行。prompt 版本为 \texttt{pair-joint-2026-11-12.1}，每对 180 条请求全部成功，0 invalid、0 API error，成本 0.153593/0.071803 美元。

防泄漏审计覆盖 180/180 对：wrapper 不含 \texttt{pair\_id}、\texttt{idx}、\texttt{commit\_id}、\texttt{filepath}、\texttt{func\_name}、\texttt{project}、CVE/CWE 列表或 ground-truth 字段名；0 个 slot leak；\texttt{requests.jsonl} 不含 \texttt{Authorization} 头，manifest 记录 \texttt{api\_key\_value\_stored=false}。唯一残余风险是 1/180 对（\texttt{18396:18397}）的两侧相同注释含 \texttt{CVE-2017-17969}；该注释不指示哪一侧是漏洞版本，但无法在不改数据的前提下清除。


% Auto-generated by report/scripts/build_extended_tables.py. Do not edit.

\begin{table*}[t]
\caption{Pair-joint 强制二选一结果（180 对；每对一次请求）}
\label{tab:pj}
\centering
\scriptsize
\setlength{\tabcolsep}{3.2pt}
\begin{tabular}{@{}lrrrrrrrrrr@{}}
\toprule
模型 & 正确 & DA [95\% CI] & 判反 & DA(A) & DA(B) & $\Delta$(A$-$B) & Fisher $p$ & 证据合规 & 定位命中 & USD\\
\midrule
DeepSeek & 138/180 & 76.7 [70.0, 82.2] & 42 & 92.2 & 61.1 & 31.1 & \ensuremath{1.00\times10^{-6}} & 98.3 & 93.0 & 0.153593 \\
MiMo & 156/180 & 86.7 [80.9, 90.9] & 24 & 88.9 & 84.4 & 4.4 & 0.5114 & 97.8 & 95.9 & 0.071803 \\
\bottomrule
\end{tabular}
\end{table*}


\begin{table*}[t]
\caption{结构化证据提示与 direct 基线（SecVulEval 160 条；valid-only 指标与 strict 全样本）}
\label{tab:ev}
\centering
\scriptsize
\setlength{\tabcolsep}{3.5pt}
\begin{tabular}{@{}llrrrrrrrrr@{}}
\toprule
模型 & 提示 & 有效 & P & R & F1+ & Macro-F1 & Strict & 预测正类率 & USD & p95 (s)\\
\midrule
DeepSeek & direct & 160/160 & 0.826 & 0.297 & 0.437 & 0.383 & 0.388 & 0.287 & 0.066988 & 10.83 \\
DeepSeek & evidence & 159/160 & 0.817 & 0.386 & 0.524 & 0.422 & 0.438 & 0.377 & 0.085622 & 12.06 \\
MiMo & direct & 150/160 & 0.806 & 0.242 & 0.372 & 0.346 & 0.325 & 0.240 & 0.024397 & 12.00 \\
MiMo & evidence & 158/160 & 0.810 & 0.370 & 0.508 & 0.407 & 0.419 & 0.367 & 0.035933 & 10.87 \\
\bottomrule
\end{tabular}
\end{table*}


\begin{table*}[t]
\caption{证据提示相对 direct 的配对检验（共同 valid；2 模型 $\times$ 2 终点 Holm）}
\label{tab:ev-cmp}
\centering
\scriptsize
\setlength{\tabcolsep}{4pt}
\begin{tabular}{@{}llrrrr@{}}
\toprule
模型 & 终点 & $\Delta$ (pp) & 不一致（新/旧） & 原始 $p$ & Holm $p$\\
\midrule
DeepSeek & 召回 & 9.4 & 1/13 & 0.0018 & 0.0055 \\
DeepSeek & FPR & 9.4 & 3/0 & 0.2500 & 0.2500 \\
MiMo & 召回 & 10.9 & 0/13 & \ensuremath{2.44\times10^{-4}} & 0.0010 \\
MiMo & FPR & 13.8 & 4/0 & 0.1250 & 0.2500 \\
\bottomrule
\end{tabular}
\end{table*}



\subsection{联合提示的结果与任务变化}
表~\ref{tab:pj} 给出主结果。DeepSeek 正确 138/180，DA 76.7\%（Wilson 95\% 区间 70.0--82.2）；MiMo 正确 156/180，DA 86.7\%（80.9--90.9）；swapped 分别为 42/24，both-vulnerable 与 both-safe 都是 0。后者不是模型“学会了强制二选一”，而是 prompt 明确要求只返回 A 或 B；它是任务约束的直接结果，不能解释为开放任务中的行为。

与分侧调用相比，DeepSeek direct 的 DA 从 11.7\% 提到 76.7\%（$\Delta=+65.0$ 个百分点，项目聚类 95\% CI $[58.6,77.9]$，McNemar $p=3.6\times10^{-34}$）；MiMo direct 从 7.8\% 提到 86.7\%（$\Delta=+78.9$，CI $[74.1,87.6]$，$p=3.6\times10^{-43}$）。但这不是纯粹的“跨版本推理”增益：分侧调用的 both-safe 为 113/108 对，联合 prompt 把它压到 0；任务从“两次自由判定”变成了“一次强制选择”。DA 的提高同时包含“看到另一版本”“被禁止输出 both/neither”“必须引用证据行”三个因素，当前实验没有消融，不能把 +65/+79 全部归因于模型理解了修复差异。

跨模型比较：MiMo 比 DeepSeek 高 10.0 个百分点（138 vs.\ 156，McNemar $p=0.0009$，项目聚类 CI $[-14.6,-4.6]$）。这个差异是真实的（在同一 180 对、同一 prompt、同一 A/B 映射上），但不能直接外推为“MiMo 的漏洞检测能力更强”，因为两者都有长度信号，且 DeepSeek 有明显位置偏差；下一节给出这两个混杂。

\subsection{位置偏差}
表~\ref{tab:pj} 的 DA(A)/DA(B) 与 Fisher 检验显示：DeepSeek 在漏洞版本位于 A 时 DA 为 92.2\%，位于 B 时只有 61.1\%，差值 31.1 个百分点（Fisher $p=1.0\times10^{-6}$）；它选择 A 的次数为 118/180（65.6\%），二项检验 $p=3.6\times10^{-5}$。MiMo 的 A/B 差值为 4.4 个百分点（88.9\% vs.\ 84.4\%，$p=0.511$），选择 A 为 94/180（52.2\%，$p=0.602$），位置偏差小且不显著。

A/B 映射由 HMAC 随机产生，位置与标签在期望上独立；分侧 direct 基线在同样的 A/B 分组上 DA 为 13.3\% vs.\ 10.0\%（DeepSeek）、10.0\% vs.\ 5.6\%（MiMo），两组难度接近，因此 DeepSeek 的 31.1 个百分点不太可能由“B 组恰好更难”解释。它更像 first-position bias：模型倾向于选择先出现的版本。当前只做了 90/90 平衡，没有做 A/B 双顺序重复，因此不能区分“位置偏好”和“该位置代码更易判”的细微交互。

\subsection{长度混淆}
修复通常会增加检查或代码，所以漏洞版本往往更短。用一个完全不看语义的规则——“选字符数较短的版本，等长时选 A”——在 180 对上正确 134 次，DA 74.4\%；按行数则为 133/180（73.9\%）。DeepSeek 的 76.7\% 与字符启发式在统计上不可区分（McNemar $b=25,c=21$，$p=0.6587$）；MiMo 的 86.7\% 显著高于启发式（$b=29,c=7$，$p=0.0003$）。

分长度层看：DeepSeek 在“漏洞更短”的 130 对上 109/130（83.8\%），在“漏洞更长”的 44 对上 25/44（56.8\%），等长 4/6；MiMo 为 123/130（94.6\%）、28/44（63.6\%）、5/6。DeepSeek 的净增益主要来自“漏洞更长”的 44 对（56.8\% vs.\ 启发式按定义的 0\%），但在“漏洞更短”的 130 对上反而低于启发式；MiMo 在两个分层都高于启发式，说明它有超出长度的信号，但总体仍受长度分布影响。该诊断是测试后加入的，不是预注册基线；没有做长度匹配子集或 A/B 反向重复，因此只能作为混杂提示，不能作为因果分解。

\subsection{结构化证据提示的协议}
另一个方向是不提供另一版本，而要求模型给出可审计的判定依据。新 prompt 在一次响应中返回六项：\texttt{is\_vulnerable}、\texttt{confidence}、\texttt{evidence\_lines}、\texttt{mechanism}、\texttt{counterfactual\_fix}、\texttt{counterfactual\_confirms}。\texttt{evidence\_lines} 是提示中打印的 1-based 行号，必须落在 \texttt{FUNCTION\_SOURCE} 内；\texttt{mechanism} 说明缺陷机制与数据流；\texttt{counterfactual\_fix} 给出最小修复；\texttt{counterfactual\_confirms} 是模型自检，不是执行验证。标签一致性规则强制：判漏洞时至少一条证据行、机制与修复不能为空；判安全时证据行为空、机制与修复为 \texttt{none}、自检为 false。

输入隔离是结构性的：构造器只接收 \texttt{func\_body}，项目名、路径、函数名、commit、CVE/CWE、\texttt{changed\_lines}、\texttt{changed\_statements}、\texttt{fixed\_func\_idx}、标签和 \texttt{context} 都不进入 prompt；用户消息的 header/footer 之间必须逐字符等于带行号源码，去掉行号前缀后必须逐行重建原函数。两个真实 run 的 160/160 条请求体 SHA-256 与 dry-run 逐条一致。DeepSeek 159/160 valid、1 invalid、0 API error；MiMo 158/160 valid、2 invalid、0 API error。

\subsection{召回上升，但操作点也移动}
表~\ref{tab:ev} 给出全样本 valid-only 指标。DeepSeek direct 的召回 29.1\%、FPR 25.0\%、macro-F1 0.383；evidence 把召回提到 38.6\%、FPR 提到 34.4\%、macro-F1 0.422、strict accuracy 从 38.8\% 提到 43.8\%。MiMo direct 的召回 24.2\%、FPR 23.3\%、macro-F1 0.346；evidence 把召回提到 37.0\%、FPR 提到 35.5\%、macro-F1 0.407、strict accuracy 从 32.5\% 提到 41.9\%。

表~\ref{tab:ev-cmp} 的配对检验（共同 valid）给出更严格的口径：DeepSeek 召回 $+9.4$ 个百分点（新增 13 个 TP、丢失 1 个 TP；原始 $p=0.0018$，Holm $p=0.0055$），FPR $+9.4$（新增 3 个 FP、修复 0 个；原始 $p=0.25$，Holm $0.25$）；MiMo 召回 $+10.9$（新增 13、丢失 0；原始 $p=2.4\times10^{-4}$，Holm $0.0010$），FPR $+13.8$（新增 4、修复 0；原始 $p=0.125$，Holm $0.25$）。召回提升在两个模型上都通过 Holm；FPR 的点估计上升，但 32/29 个负例的功效不足以把它判为显著。因此可以说“召回提高有统计支持”，不能说“没有增加假阳性”，更不能说“判别力提高”。

操作点分析支持这一解释：共同 valid 上预测正类率从 28.3\% 升到 37.7\%（DeepSeek）、23.7\% 升到 35.1\%（MiMo）；$\Delta$balanced accuracy 为 $+0.0004$/$-0.0143$，$\Delta$MCC 为 $-0.0020$/$-0.0276$。F1 和 strict accuracy 的上升主要来自“更敢报正例”，而不是排序能力变好。MiMo 的 strict 提升还混入格式合规恢复：direct 有 10 条 invalid（8 正/2 负），evidence 只有 2 条；共同 valid 的 accuracy 提升是 $+6.1$ 个百分点，小于 strict 的 $+9.4$。

\subsection{证据行审计：可审计不等于可定位}
所有 valid 正例的证据行都在范围内、非空、无重复（DeepSeek 60/60，MiMo 58/58 grounded）。与 ground-truth 修复行（\texttt{changed\_lines} 映射回函数体）的重合：DeepSeek 的 exact hit 为 48.8\%，同函数同数量随机引用基线为 34.5\%；MiMo 为 45.5\% vs.\ 31.6\%。$\pm3$ 重合为 62.8\% vs.\ 62.5\%（DeepSeek）、63.6\% vs.\ 59.7\%（MiMo），几乎与随机基线相同。引用行的 exact 精确率只有 21.9\%/25.6\%，修复行的 exact 召回只有 16.3\%/14.6\%。

反事实自检的确认率为 100\%/98.3\%，但在假阳性上也是 100\%；它没有提供降 FP 的信号。事后“只保留 grounded 且已确认正例”的敏感性分析几乎不改变指标，说明该自检更像一致性检查，而不是独立 oracle。结论：证据提示的贡献是“让模型给出可审计的结构”，不是“定位到了缺陷行”；把 \texttt{evidence\_lines} 当作人工复核排序信号，会被随机基线解释掉大部分。

\section{扩展评估 III：Juliet CWE-401 分层基准}
\label{sec:ext-juliet}

Tier A 的真实项目函数片段无法回答“同一 CWE-401 用例上，静态/动态/LLM/混合各自如何”。本节使用 NIST Juliet v1.3 的 CWE-401 子集构建 Tier B，单独报告，不与 Tier A 合并。Juliet 的标签由宏结构直接给出，可编译可运行；代价是合成、单 CWE、模板化，且不同家族之间近重复程度高。

\subsection{数据构建与泄漏控制}
数据来自 NIST SARD 发布的 Juliet Test Suite for C/C++ v1.3，CWE-401 Memory Leak；归档 SHA-256 为 \texttt{ada9d7e1...f38eb}，提取 1{,}228 个 \texttt{.c}、1{,}482 个 \texttt{.cpp}、111 个 \texttt{.h}。从 1{,}658 个候选组中按 42 个功能家族确定性选择 100 组，每组一个 bad、一个 good，形成 200 个构建单元；dev 60（30 组、13 家族）、test 140（70 组、29 家族）；C 88、C++ 112。200/200 编译通过，200/200 宏门控与符号门控通过；good 100/100 默认运行无泄漏，bad 78/100 默认运行观察到泄漏，22/100 为条件性/可达性情形（15 个 \texttt{malloc\_realloc} 只在 \texttt{realloc} 失败时泄漏，7 个 \texttt{\_68} 的指针在退出时仍被全局变量引用）。

dev/test 按近重复簇隔离：去掉注释、统一标识符/字面量后取 5-gram shingle，Jaccard 阈值 0.85，整簇分配；跨 split 2{,}100 对中 $\ge0.85$ 为 0 对，最大 0.8231。这个划分不是 i.i.d. 随机划分：dev 13 个家族、test 29 个家族，分布有意不同，适合“dev 调参、test 测泛化”，不适合把两侧当同分布做显著性检验。


% Auto-generated by report/scripts/build_extended_tables.py. Do not edit.

\begin{table*}[t]
\caption{Juliet CWE-401 静态/动态结果（test 140：70 bad / 70 good；valid-only macro-F1）}
\label{tab:jb-static}
\centering
\scriptsize
\setlength{\tabcolsep}{3.2pt}
\begin{tabular}{@{}lrrrrrrrrr@{}}
\toprule
方法 & 有效 & 覆盖 & TP/FP/TN/FN & P & R & Spec & Macro-F1 [95\% CI] & MCC & Strict\\
\midrule
Clang SA 22.1.8 & 136/140 & 0.971 & 53/0/68/15 & 1.000 & 0.779 & 1.000 & 0.888 [0.818, 0.948] & 0.799 & 0.864 \\
Cppcheck 2.21.0 & 140/140 & 1.000 & 22/4/66/48 & 0.846 & 0.314 & 0.943 & 0.588 [0.497, 0.671] & 0.331 & 0.629 \\
ASan+LSan & 140/140 & 1.000 & 58/0/70/12 & 1.000 & 0.829 & 1.000 & 0.914 [0.850, 0.966] & 0.841 & 0.914 \\
ASan+LSan+Valgrind & 140/140 & 1.000 & 58/0/70/12 & 1.000 & 0.829 & 1.000 & 0.914 [0.850, 0.966] & 0.841 & 0.914 \\
Valgrind 3.27.1 & 140/140 & 1.000 & 58/0/70/12 & 1.000 & 0.829 & 1.000 & 0.914 [0.850, 0.966] & 0.841 & 0.914 \\
\bottomrule
\end{tabular}
\end{table*}


\begin{table*}[t]
\caption{动态保守视图：clean 弃权（test 140）}
\label{tab:jb-dyncons}
\centering
\scriptsize
\setlength{\tabcolsep}{3.5pt}
\begin{tabular}{@{}lrrrrrrrr@{}}
\toprule
方法 & 有效 & 覆盖 & TP & FP & FN & 弃权 & 召回 & 精确率\\
\midrule
ASan+LSan & 58/140 & 0.414 & 58 & 0 & 0 & 82 & 1.000 & 1.000 \\
ASan+LSan+Valgrind & 58/140 & 0.414 & 58 & 0 & 0 & 82 & 1.000 & 1.000 \\
Valgrind 3.27.1 & 58/140 & 0.414 & 58 & 0 & 0 & 82 & 1.000 & 1.000 \\
\bottomrule
\end{tabular}
\end{table*}


\begin{table*}[t]
\caption{Juliet CWE-401 LLM 直接检测（MiMo，test 140；Xiaomi 官方端点 fallback；修复前运行仅作审计）}
\label{tab:jb-llm}
\centering
\scriptsize
\setlength{\tabcolsep}{4pt}
\begin{tabular}{@{}lrrrrrrrrr@{}}
\toprule
运行 & 有效 & TP & FP & TN & FN & Acc (\%) & F1 & MCC & 参考 USD\\
\midrule
修复后正式 & 140/140 & 68 & 5 & 65 & 2 & 95.0 & 0.951 & 0.901 & 0.003180 \\
修复后复现 & 140/140 & 68 & 5 & 65 & 2 & 95.0 & 0.951 & 0.901 & 0.003081 \\
修复前（仅审计） & 140/140 & 68 & 4 & 66 & 2 & 95.7 & 0.958 & 0.915 & 0.004426 \\
\bottomrule
\end{tabular}
\end{table*}


\begin{table*}[t]
\caption{Juliet 证据混合 v2（泄漏修复后）：2 模型 $\times$ 5 方法（test 140；strict 全样本；1{,}400/1{,}400 请求成功）}
\label{tab:hybrid}
\centering
\scriptsize
\setlength{\tabcolsep}{3.2pt}
\begin{tabular}{@{}llrrrrrrrr@{}}
\toprule
模型 & 方法 & 正确/140 & Acc (\%) & P & R & F1 & Spec & MCC & 列表 USD\\
\midrule
DeepSeek & direct & 135/140 & 96.43 & 0.933 & 1.000 & 0.966 & 0.929 & 0.931 & 0.010185 \\
DeepSeek & static & 135/140 & 96.43 & 0.933 & 1.000 & 0.966 & 0.929 & 0.931 & 0.011730 \\
DeepSeek & dynamic & 136/140 & 97.14 & 0.946 & 1.000 & 0.972 & 0.943 & 0.944 & 0.011471 \\
DeepSeek & combined & 136/140 & 97.14 & 0.946 & 1.000 & 0.972 & 0.943 & 0.944 & 0.011891 \\
DeepSeek & shuffled & 134/140 & 95.71 & 0.921 & 1.000 & 0.959 & 0.914 & 0.918 & 0.010860 \\
MiMo & direct & 136/140 & 97.14 & 0.946 & 1.000 & 0.972 & 0.943 & 0.944 & 0.004718 \\
MiMo & static & 135/140 & 96.43 & 0.933 & 1.000 & 0.966 & 0.929 & 0.931 & 0.006287 \\
MiMo & dynamic & 133/140 & 95.00 & 0.932 & 0.971 & 0.951 & 0.929 & 0.901 & 0.005735 \\
MiMo & combined & 136/140 & 97.14 & 0.946 & 1.000 & 0.972 & 0.943 & 0.944 & 0.004813 \\
MiMo & shuffled & 133/140 & 95.00 & 0.920 & 0.986 & 0.952 & 0.914 & 0.902 & 0.004645 \\
\bottomrule
\end{tabular}
\end{table*}


\begin{table*}[t]
\caption{Juliet 混合的配对比较（v2；pp；新对/旧对为方法对基线）}
\label{tab:hybrid-cmp}
\centering
\scriptsize
\setlength{\tabcolsep}{4pt}
\begin{tabular}{@{}llrrr@{}}
\toprule
模型 & 对比 & $\Delta$ pp [95\% CI] & 新对/旧对 & McNemar $p$\\
\midrule
DeepSeek & direct$\to$dynamic & 0.71 [0.00, 2.14] & 1/0 & 1.0000 \\
DeepSeek & direct$\to$combined & 0.71 [0.00, 2.14] & 1/0 & 1.0000 \\
DeepSeek & direct$\to$shuffled & -0.71 [-2.14, 0.00] & 0/1 & 1.0000 \\
DeepSeek & combined$\to$shuffled & -1.43 [-3.57, 0.00] & 0/2 & 0.5000 \\
MiMo & direct$\to$static & -0.71 [-2.14, 0.00] & 0/1 & 1.0000 \\
MiMo & direct$\to$dynamic & -2.14 [-5.00, 0.00] & 0/3 & 0.2500 \\
MiMo & direct$\to$combined & 0.00 [-2.14, 2.14] & 1/1 & 1.0000 \\
MiMo & direct$\to$shuffled & -2.14 [-5.00, 0.00] & 0/3 & 0.2500 \\
MiMo & combined$\to$shuffled & -2.14 [-5.00, 0.00] & 0/3 & 0.2500 \\
\bottomrule
\end{tabular}
\end{table*}


\begin{table*}[t]
\caption{Juliet 混合 v1$\to$v2 审计（v1 为清洗前；v2 为主结果；$\Delta$ = v2$-$v1）}
\label{tab:hybrid-audit}
\centering
\scriptsize
\setlength{\tabcolsep}{4pt}
\begin{tabular}{@{}llrrrrr@{}}
\toprule
模型 & 方法 & v1 正确/140 & v2 正确/140 & $\Delta$ pp [95\% CI] & 翻转（对→错/错→对） & McNemar $p$\\
\midrule
DeepSeek & direct & 135/140 & 135/140 & +0.00 [+0.00, +0.00] & 0/0 & 1.0000 \\
DeepSeek & static & 135/140 & 135/140 & +0.00 [+0.00, +0.00] & 0/0 & 1.0000 \\
DeepSeek & dynamic & 139/140 & 136/140 & -2.14 [-5.00, +0.00] & 3/0 & 0.2500 \\
DeepSeek & combined & 140/140 & 136/140 & -2.86 [-5.71, -0.71] & 4/0 & 0.1250 \\
DeepSeek & shuffled & 135/140 & 134/140 & -0.71 [-3.57, +1.43] & 2/1 & 1.0000 \\
MiMo & direct & 135/140 & 136/140 & +0.71 [+0.00, +2.14] & 0/1 & 1.0000 \\
MiMo & static & 136/140 & 135/140 & -0.71 [-2.14, +0.00] & 1/0 & 1.0000 \\
MiMo & dynamic & 135/140 & 133/140 & -1.43 [-3.57, +0.00] & 2/0 & 0.5000 \\
MiMo & combined & 138/140 & 136/140 & -1.43 [-3.57, +0.00] & 2/0 & 0.5000 \\
MiMo & shuffled & 135/140 & 133/140 & -1.43 [-3.57, +0.00] & 2/0 & 0.5000 \\
\bottomrule
\end{tabular}
\end{table*}



\subsection{静态与动态结果}
表~\ref{tab:jb-static} 给出 test 140 的主结果。Clang Static Analyzer 22.1.8 有效 136/140（4 个 \texttt{\_82} 组因共享无 include-guard 头而 \texttt{compile\_error}），macro-F1 0.888（95\% CI 0.818--0.948），MCC 0.799，strict 0.864；它没有 FP（精确率 1.000），但漏掉 15 个 bad（召回 0.779）。Cppcheck 2.21.0 全部有效，macro-F1 0.588（0.497--0.671），MCC 0.331，strict 0.629；精确率 0.846、召回 0.314，4 个 FP、48 个 FN。ASan+LSan 与 Valgrind 在 test 140 上给出完全相同的 operational 结果：58 TP、0 FP、70 TN、12 FN，macro-F1 0.914（0.850--0.966），MCC 0.841，strict 0.914。ASan+Valgrind 的“任一报泄漏”合并视图与二者相同。

静态的 4 个 \texttt{compile\_error} 保留在 strict 分母中，但从未写成 $y_{\mathrm{pred}}=0$；valid-only macro-F1 排除它们，strict 把它们算错。这个区别很重要：如果把编译失败当作“无泄漏”，Clang 的 strict 会从 0.864 变成 0.889，看起来更好，但掩盖了工具在 4 个单元上根本没有给出预测。Cppcheck 的 coverage 是 1.000，但召回只有 0.314；高 coverage 不等于高检测率。

\subsection{动态保守视图与 runtime-supported 分层}
表~\ref{tab:jb-dyncons} 把动态 clean 全部弃权。test 140 中只有 58 个单元有可观测泄漏，coverage 0.414；82 个 clean 预测被弃权，负类因此离开 valid 集，macro-F1 未定义。在可观测的 58 个单元上，召回与精确率都是 1.000，但这只是“观察到泄漏的都报对了”，不是安全分类器的性能。动态 run 没有行/分支覆盖率 instrumentation；clean 只表示“执行路径上未观察到 definite/indirect leak”，不是 TN，也不是安全证明。

runtime-supported 视图（test 128）排除 12 个“标签 bad 但 LSan 未观察到泄漏”的单元。ASan/Valgrind 在该子群上达到 1.000，但该子群本身由 LSan preflight 定义，存在循环性；Clang 在该子群上 0.917、Cppcheck 0.564，静态与 Valgrind 的数字更有信息量。报告同时给 operational 与 conservative 两个视图，不允许把二者混用：前者把 clean 当 TN，后者把 clean 弃权；前者适合“工具在默认执行路径上能抓到多少”，后者适合“没有覆盖率证据时能声称多少”。

\subsection{Juliet 上的 LLM 直接检测}
表~\ref{tab:jb-llm} 给出 MiMo 在 test 140 上的三次真实运行。泄漏修复后的正式与复现运行都是 140/140 有效，TP 68、FP 5、TN 65、FN 2，Accuracy 95.00\%，F1 0.951，MCC 0.901；两次运行的混淆矩阵完全相同，说明该结果在该端点上可重复。修复前的运行保留作审计：TP 68、FP 4、TN 66、FN 2，Accuracy 95.71\%，F1 0.958；修复后只多 1 个 FP，方向与清洗掉 6 个 good 单元的标签字面量一致，但只有 1 个单元变化，不能把差值单独归因于清洗。主结果使用修复后的 95.0\%，修复前的 95.7\% 不进入结论。分语言：C 58 单元 Accuracy 93.10\%，C++ 82 单元 96.34\%。评分按实际 \texttt{label} 字段，而不是按模型理由文本；7 个错误为 5 FP + 2 FN。

必须标明端点：该运行使用 Xiaomi 官方 OpenAI 兼容端点（\url{https://api.xiaomimimo.com/v1}）的 \texttt{mimo-v2.6-flash}，不是 OpenCode Go；相同模型 ID 不等于相同服务栈，延迟、缓存、成本和可能的系统提示注入都可能不同。DeepSeek 在 Juliet 上没有结果：OpenCode Go 返回 401 \texttt{api\_key\_expired}，DeepSeek 官方返回 402 \texttt{Insufficient Balance}，qwen 网关返回 403 \texttt{AccessDenied.Unpurchased}；没有用其他模型或 prompt 冒充。参考成本为正式 0.003180 美元、复现 0.003081 美元（按 MiMo 发布的 0.14/0.28/0.0028 计算），不是 Xiaomi 实际账单；Xiaomi 端点不返回 provider list/allowance，故该成本只有 configured/normalized 一列。

登记口径也需要说明：\path{benchmark/results/tier_b_summary.json} 仍把 B4--B8 标为 \texttt{blocked\_gate}，理由是 B9 消毒/标签审计未完成、没有登记为统一 schema 的 Juliet LLM 请求；而 \path{juliet_llm/} 是独立子工作流，已有上述真实运行。本文把它作为“未纳入统一 leaderboard 的探索性证据”报告，不把它写成 Tier B B4/B5 的正式结果，也不把 95.0\% 与 Tier A 的 0.383 macro-F1 直接比较。

\subsection{证据混合与安慰剂}
\path{hybrid/} 在同样的 test 140 上运行 2 个模型 $\times$ 5 种方法：direct、static、dynamic、combined、shuffled。direct 与 \path{juliet_llm} 的 prompt 逐字节一致（140/140）；static 证据来自真实 Clang/Cppcheck 行，dynamic 证据来自真实 ASan/LSan 与 Valgrind，combined 是二者拼接，shuffled 用另一个单元（不同 group、非自身）的 combined 证据作为安慰剂，seed 20261112。证据文本清除绝对路径、文件名、行号、CWE 标识、\texttt{OMIT*} 宏和 bad/good 标识符；inconclusive 明确写“不是 clean”。首轮 v1 完成后，发布审计发现 6 个 good 单元仍含字面量 \texttt{"Benign, fixed string"}；共享 cleaner 修复后，用独立输入包 \path{hybrid/inputs_v2} 与 prompt v2 重跑全部 10 个 run、1{,}400 次请求（v2）。预检 PASS，10/10 run 均为 140/140、0 invalid、0 API error。本小节以 v2 为主结果，v1 冻结保留、只作审计对照。

表~\ref{tab:hybrid} 给出 v2 的 strict 结果。DeepSeek：direct 135/140（96.43\%）、static 135、dynamic 136（97.14\%）、combined 136（97.14\%）、shuffled 134（95.71\%）；MiMo：direct 136（97.14\%）、static 135、dynamic 133（95.00\%）、combined 136（97.14\%）、shuffled 133（95.00\%）。表~\ref{tab:hybrid-cmp} 的配对检验显示，v2 中没有任何方法相对 direct 显著：DeepSeek direct$\to$dynamic 与 direct$\to$combined 都是 $+0.71$ 个百分点（1 对 0，McNemar $p=1.0$，CI $[0.00,2.14]$），direct$\to$shuffled 为 $-0.71$（0 对 1，$p=1.0$）；MiMo direct$\to$combined 为 $0.00$（1 对 1，$p=1.0$），direct$\to$shuffled 为 $-2.14$（0 对 3，$p=0.25$）。v2 的最大单元格差异只有 2.14 个百分点，且方向并不一致；证据混合的收益在修复后的数据上更弱。

表~\ref{tab:hybrid-audit} 把 v1 与 v2 并列，用于量化清洗的影响，而不是给 v1 排名。DeepSeek combined 从 140/140（100\%）降到 136/140（97.14\%），差 $-2.86$ 个百分点（4 对 0，$p=0.125$，CI $[-5.71,-0.71]$）；DeepSeek dynamic 从 139 降到 136。6 个被清洗的单元共引起 17 次标签翻转：16 次由对变错、1 次由错变对（DeepSeek shuffled 的 \texttt{malloc\_realloc\_int64\_t\_09\_good}）；134 个未变单元只有 1 次端点噪声翻转（MiMo direct 的 1 个 bad）。五方法池化准确率从 97.36\% 降到 96.36\%。因此 v1 的 100\% 是标签捷径下的上界，只能作为“清洗前后差多少”的审计案例；主结论使用 v2，DeepSeek combined 为 97.14\%，与 direct 的差异在 1 个单元、精确 McNemar $p=1.0$，不能声称组合证据带来稳定提升。

v2 的安慰剂没有聚合收益：DeepSeek shuffled 134/140 比 direct 少 1 个，MiMo 133/140 少 3 个；相对 direct 只有 1/3 个预测发生翻转，且全部由对变错。按供体标签统计，DeepSeek 在供体 bad 时判 vulnerable 的比例为 57.1\%、供体 good 为 51.4\%；MiMo 为 55.7\%/51.4\%。错配证据仍会按供体方向轻微推动边界样本，但没有带来净收益。这说明 combined 的 $+1/0$ 个单元不能简单归因于“多了一段证据文本”：同样长度、同样格式的错配证据没有收益甚至略差。但受样本量与单一 derangement 限制，这仍只是方向性证据；另有 14 个单元的供体证据与自身完全相同（no-op collision），在这些单元上安慰剂无法改变 prompt。

v2 的错误结构仍集中在 \texttt{\_22} good：DeepSeek direct 的 5 个 FP 全是 \texttt{\_22} good（\texttt{int\_malloc\_22}、\texttt{new\_array\_char\_22}、\texttt{new\_array\_int\_22}、\texttt{new\_int\_22}、\texttt{twoIntsStruct\_realloc\_22}），dynamic 与 combined 各纠正 1 个，仍有 4 个；MiMo direct 与 combined 的 4 个 FP 也全是 \texttt{\_22} good，dynamic 额外产生 2 个 FN（\texttt{malloc\_realloc\_int\_01\_\_bad}、\texttt{malloc\_realloc\_int\_18\_\_bad}）。\texttt{\_22} 变体通过全局变量在多个 sink 之间传递控制流，good 版本在“另一条”分支释放内存；direct 容易只看到当前分支未释放而误报。动态 clean 与组合证据只能纠正部分单元，天花板效应明显：direct 已 135/140，证据最多只能再修复 5 个。

同样需要登记口径：\path{benchmark/results/tier_b_summary.json} 把 B6--B8 标为 \texttt{blocked\_gate}，而 \path{hybrid/} 是独立工作流，已有 v2 的 1{,}400 次真实请求。本文把 hybrid 作为“未纳入统一 leaderboard 的探索性证据”，不写成 B6--B8 的正式结果；它的价值是给出安慰剂对照和清洗审计，而不是给 Tier B 增加一行排名。

同一 MiMo 模型在 \path{juliet_llm/} 的 Xiaomi 端点上（修复后）direct 为 133/140（95.00\%），在 \path{hybrid/} 的 OpenCode Go 上 direct 为 136/140（97.14\%）；差 3 个单元。两个工作流的 prompt 逐字节一致，且输入都已使用修复后的清洗版本，但端点、缓存和可能的系统提示注入不同，且各只有单次运行，不能把这 3 个单元当作模型差异。

\subsection{跨层不比较}
Tier B 的 Clang 0.888 不能与 Tier A 的 LLM 0.383 比较：前者是 140 个合成 CWE-401 单元、70/70 平衡、标签由宏结构给出；后者是 160 个真实项目函数片段、128/32 不平衡、负例只是“源数据未标注漏洞”。动态 0.914 也不能与 LLM 的 0.950 比较：动态的 0.914 把 clean 当 TN，而 clean 没有覆盖率证据；LLM 的 0.950 是在同一个合成集上、用 sanitized 源码、单次运行得到。跨层数字的唯一合法用途是分别回答“真实项目函数上的提示策略如何”和“合成 CWE-401 上的工具/模型行为如何”，不是产生一个“综合分”。

\section{扩展评估 IV：修复的编译--LSan--行为闭环}
\label{sec:ext-repair}

旧报告的修复评估只有结构代理和 LLM 辅助盲评，所有候选的 \texttt{compile\_status} 都是 \texttt{NOT\_RUN}。扩展部分在 Juliet CWE-401 的 70 个 test bad 单元上补上了真实执行 oracle：模型输出必须能编译、在 LSan 下不再报泄漏、并且保持可观测行为，才算严格通过。

\subsection{协议}
输入是评测端净化的“bad-only”源码：删除全部注释（包括 \texttt{POTENTIAL FLAW}、\texttt{FIX}、\texttt{BadSink} 等提示）、用预处理器求值删除 \texttt{OMITGOOD} 区域、删除 \texttt{INCLUDEMAIN} 区域、删除只属于 good 路径的全局声明；评测端另写 probe \texttt{main()} 只调用 bad 入口。模型永远看不到 good 实现，也看不到“哪个文件是 good 参考”。自动审计检查 \texttt{FIX}、\texttt{POTENTIAL FLAW}、\texttt{GoodSink}、\texttt{BadSink}、\texttt{good*} 标识符、\texttt{good()} 调用、块注释和行注释；70/70 单元 0 命中。\texttt{check\_baselines.py} 另外证明净化没有改变语义：70/70 净化输入的编译、泄漏观测和行为行与原始 bad 完全一致。

编译使用 \texttt{gcc -std=gnu11} / \texttt{g++ -std=gnu++11}，\texttt{-O0 -g -fsanitize=address -fno-omit-frame-pointer}，链接 Juliet 的 \texttt{io.o} 与 \texttt{std\_thread.o}。运行时 \path{ASAN_OPTIONS=detect_leaks=1:abort_on_error=0:exitcode=23}，\path{LSAN_OPTIONS=exitcode=23:use_stacks=0:use_registers=0}；不把栈/寄存器里的悬垂指针当作 LSan 根，避免“修完但栈上残留指针”被误判为无泄漏；全局变量仍作为根保留，因为 \texttt{\_68} 变体的唯一指针在全局变量里。行为判定用 probe 的 stdout 数据行与原始 bad 逐行比较，要求完全相等、exit code 0、无其他 ASan 错误/超时。严格通过定义为 \texttt{compile\_pass} \&\& \texttt{leak\_removed} \&\& \texttt{behavior\_preserved} \&\& \texttt{baseline\_leak\_observed}；baseline 未观察到泄漏的单元一律不计 strict pass。每个单元先请求一次，失败后最多一次反馈；反馈只含编译器 stderr 尾部、LSan summary 和可观测输出的期望/实际值，不含 good 源码。


% Auto-generated by report/scripts/build_extended_tables.py. Do not edit.

\begin{table*}[t]
\caption{Juliet 修复闭环：编译 + LSan + 行为差分（70 个 test bad 单元）}
\label{tab:repair-close}
\centering
\scriptsize
\setlength{\tabcolsep}{4pt}
\begin{tabular}{@{}lrrrrrrrr@{}}
\toprule
模型 & 编译/70 & LSan 干净/70 & 行为保持/70 & 严格通过/58 & 首次/58 & 迭代单元 & 调用 & USD\\
\midrule
DeepSeek & 70/70 & 70/70 & 70/70 & 58/58 & 58/58 & 0 & 70 & 0.026990 \\
MiMo & 70/70 & 70/70 & 70/70 & 58/58 & 57/58 & 2 & 72 & 0.005841 \\
\bottomrule
\end{tabular}
\end{table*}


\subsection{结果}
表~\ref{tab:repair-close} 给出结果。DeepSeek 70/70 编译、70/70 LSan 干净、70/70 行为保持；58 个 baseline 可观测单元上 58/58 严格通过，首次回答即 58/58，0 次反馈，70 次调用，configured 成本 0.026990 美元（normalized 0.013495）。MiMo 同样 70/70 编译、70/70 LSan 干净、70/70 行为保持；58 个可观测单元上 58/58 严格通过，但首次只有 57/58，2 个单元使用了反馈（1 次真实 double-free、1 次 schema 路径错误），72 次调用，configured 成本 0.005841 美元（normalized 相同）。两个 headline run 合计 configured 0.032831 美元（normalized 0.019336）；加上解析器开发/冒烟的全部真实 API 成本，configured all-in 为 0.082550 美元（normalized 0.052549），而不是旧 README 的 0.077516。差额 0.005034 来自 \texttt{parser\_dev/mimo\_test70\_v1}。

12 个单元不计入 strict 分母：7 个 \texttt{malloc\_realloc} 变体只有在 \texttt{realloc} 失败时才泄漏初始块，而默认执行路径 \texttt{realloc} 成功并随后 \texttt{free}；5 个 \texttt{\_68} 变体把唯一指针放进全局变量，LSan 视为 reachable。它们仍全部编译通过且行为保持，但不能证明泄漏被移除。没有用 \texttt{ulimit}/分配器故障注入强制 \texttt{realloc} 失败，因为 ASan 需要巨大的虚拟地址空间，\texttt{RLIMIT\_AS} 不可用；\texttt{hard\_rss\_limit\_mb} 也不能可靠地让 \texttt{realloc} 返回 NULL。因此这 7 个 latent 单元保持不可观测，报告不把它们算成“白捡的通过”。

\subsection{解释边界}
这个闭环比旧报告的结构代理强得多：它证明模型输出真的能编译、真的在 LSan 下干净、真的保持了 stdout 行为。但它仍不是“修复正确率”的完整 oracle：行为只覆盖 stdout 数据行，没有覆盖文件、网络、全局状态或 API 契约；没有故障注入，12 个 latent 单元不可判定；只做一次反馈，不是多轮搜索；Juliet 是合成单 CWE，不能外推到真实项目。两个模型在 58 个可观测单元上都 58/58，说明在这个受限任务上“编译 + LSan + 行为”门槛已经被跨过；差异主要体现在首次通过率、反馈使用和成本，而不是最终 strict 通过数。这个结果不能与旧报告的 LLM 辅助盲评 17/81 直接比较：任务、输入、评分 rubric 和 oracle 都不同，一个是“rubric 下根因已处理”，一个是“可执行闭环严格通过”。

\section{创新点}
\label{sec:innovation}

本文的创新点是评估协议与证据审计，不是新模型或新训练方法。具体包括：
\begin{enumerate}
    \item \textbf{统一契约、分层比较。} 用同一套字段、失败分类、coverage/strict、Wilson/bootstrap/McNemar 和成本向量，同时约束 Tier A 与 Tier B；两层各自设主终点与外推边界，明确拒绝跨层池化。这样既能并列阅读，又不会制造“平均 F1”。
    \item \textbf{双模型复制与配置对齐。} 把三种提示策略完整复制到 MiMo-V2.6-Flash，每模型 2{,}484 次请求，并把 MiMo 的 \texttt{reasoning\_effort=none}、\texttt{max\_completion\_tokens} 与 DeepSeek 的非思考模式语义对齐；报告明确把差异解释为“模型 + 传输/解码配置”，而不是纯模型能力。
    \item \textbf{强制二选一联合提示 + 位置/长度审计。} 每对一次请求、HMAC 平衡 A/B、强制选择、证据行引用；随后用 Fisher 检验检查位置偏差，用“选较短版本”的平凡启发式检查长度混淆。正是这个审计把 DeepSeek 的 76.7\% 从“跨版本推理”降级为“与长度启发式不可区分”。
    \item \textbf{结构化证据提示 + 随机引用基线。} 要求模型给出证据行、机制、反事实修复和自检；把证据行与 \texttt{changed\_lines} 的 exact/$\pm3$ 重合率和同函数同数量随机引用基线比较。结果显示 exact 略高于随机、$\pm3$ 接近随机，从而阻止了“证据行 = 缺陷定位”的过度主张。
    \item \textbf{安慰剂对照的证据混合。} 在 Juliet 上用 5 种方法（direct/static/dynamic/combined/shuffled）和 1{,}400 次真实请求，把“证据内容”与“多了一段同格式文本”分开；v2 中 shuffled 没有聚合收益（DeepSeek 134/140、MiMo 133/140，均不高于 direct），因此 combined 的 $+1/0$ 个单元不能简单归因于提示变长。
    \item \textbf{带可评估分母的修复闭环。} 用编译、LSan、行为差分和 baseline 可观测性四重条件定义 strict pass，把 12 个不可观测单元从分母中排除；两个模型在 58 个可观测单元上均 58/58，给出了比文本匹配强得多的执行证据。
\end{enumerate}
这些设计的共同点是：先定义“什么算失败、什么算不可评估、什么算混杂”，再报告点估计。它们的主要价值不是把数字做大，而是让负面结果、微弱结果和不可评估结果都能被解释。

\section{局限性}
\label{sec:limitations}
本文的限制主要体现在以下方面。
\begin{enumerate}
    \item \textbf{样本规模与构成。} 源子集有 200 条独立函数和 200 对版本，冻结测试集为 160 条检测函数（128 正/32 负）和 180 对配对版本，足以支持课程项目的比较，但不足以覆盖 C/C++ 漏洞的全部类型。80/20 比例是人为设定，配对集的平衡标签是配对要求的结果；两者都不能用来估计真实项目中的漏洞发生率。检测测试集只有 32 条负例，特异度的置信区间会很宽。
    \item \textbf{标签继承。} 本文不重新标注数据，无法排除源数据的标签噪声、修复关联错误或“修复版本仍不安全”的情况。配对集的联合指标以数据集标签为准，若标签本身有争议，DA 的解释力有限。
    \item \textbf{任务内隔离、跨任务重叠。} 检测任务内 40 个开发项目与 40 个测试项目没有交集，配对任务内 20 个开发项目与 60 个测试项目没有交集，开发/测试源索引交集为 0。但两个任务独立划分，2 个检测开发项目出现在配对测试（yara、znc），5 个配对开发项目出现在检测测试（Fast-DDS、OpenSC、\texttt{file}、htmldoc、libgd）；源子集还共享 26 个项目。模型可能利用项目风格或相似代码，而不是安全语义，因此不能声称完全的项目级泛化。
    \item \textbf{开发集规模有限。} 冻结划分与课程建议一致（检测 40/160；配对 20/180），但开发集只有 40 条检测函数和 20 对配对版本，固定示例只有 4 条检测 + 2 条修复，未做系统的提示/示例消融。pilot 的 36 条检测和 18 对配对也不进入主结果；因此本文不能报告“在开发集上充分调参后冻结”的完整流程，示例选择仍可能影响结果。
    \item \textbf{上下文与标签相关。} 检测测试集 context 非空层 126 条全部为漏洞，空层 34 条中 2 条漏洞、32 条非漏洞（源独立集非空 156/156、空 44 条中 4 条漏洞）；非空层是单类别，M3 的总体差异与 context 可用性高度混杂。即使按 context 是否非空分层，空层只有 34 条，样本量有限，无法完全分离“上下文语义”和“有无上下文”的影响。
    \item \textbf{源级截断。} 源子集 5 条函数正文恰好 32{,}767 字符且全部在中途结束（独立 \texttt{idx}=9991、8265、7018，配对 \texttt{idx}=8675、8676），这 5 条全部落在测试集。32{,}767 是 Excel 单元格字符上限，因此长度分布是截断后的下界。超长函数的检测和修复结果只能解释为截断后的下界；8675/8676 同时属于 CWE-401 案例对，其中 8675 的输入不完整，不能当作干净的泄漏案例。
    \item \textbf{CWE-401 与逐 CWE 分层不足。} 源子集 CWE-401 漏洞函数只有 4/360（独立 2/160，配对 2/200），测试集为 3/308；其中 3 条是 \texttt{.c}、1 条是 \texttt{.cc}，唯一的 C++ 记录 17743 同时带 CWE-415 且落在开发集，因此严格“C++ 扩展名且单一 CWE-401 标签”为 0 条。4 条源记录只来自 3 个项目，linux 占 2 条。源子集合并 57 种数字 CWE，独立集 23/39 种、配对集 18/39 种只出现 1 次，只有 6/39 和 7/39 种达到 $n\ge 10$；测试集合并 54 种，29 种为单例，10 种达到 $n\ge 10$。因此本文的主定位是全 CWE 的混合评估，CWE-401 只能作案例系列，不能估计 CWE-401 的 precision/recall/F1。Tier A 的 SecVulEval 数据中 0 条 Juliet 记录；扩展部分的 Tier B Juliet 子集另建，只在该层内报告，不能与 Tier A 的指标比较或外推。
    \item \textbf{字段泄漏与标签相关。} 测试集 \texttt{fixed\_func\_idx} 非空为有漏洞 115/128、无漏洞 0/32（源独立集 143/160 与 0/40）；CWE/CVE 在测试集 32/32 条负例中也非空（源独立集 40/40），配对修复版 180/180（源 200/200）复制漏洞版。协议从所有提示中排除这些字段，但字段分布本身带有强标签信号：任何未按白名单过滤的提示、缓存或后处理都可能让指标虚高。即使正确排除，这些相关性仍意味着“字段是否可用”与标签绑定，属于构造层面的混淆。
    \item \textbf{单一基础模型。} 本文比较的是提示策略，不是模型能力。结论不能推广到其他模型、其他版本或其他服务商；模型更新后需要重跑。
    \item \textbf{提示敏感性。} 提示措辞、示例选择、上下文格式和输出模式都可能影响结果。本文固定一套提示并公开模板，但没有做系统的提示消融；因此策略间的差异包含提示设计的偶然性。每种策略对每条记录只取一次有效判定（解析失败的重试不计为重复采样），因此无法估计服务端非确定性带来的方差。
    \item \textbf{修复没有 oracle。} 修复任务没有编译、测试、PoC 或语义等价验证，只能报告结构代理；精确匹配参考修复是文本相似度，不是正确性。第~\ref{sec:human-review}节的盲评已完成，但它是两名 LLM 辅助评审加第三 LLM 裁决，不是人类评审；所有候选的 compile\_status 为 NOT\_RUN，MIN 的二次加权 $\kappa$ 仅 0.015。因此不能声称“修复成功”或报告修复正确率；rubric 通过计数只是该 LLM 辅助 rubric 下的判断，不是可执行 oracle。这里的 NOT\_RUN 指旧报告的 SecVulEval 修复候选；扩展部分第~\ref{sec:ext-repair}节对 Juliet 的 70 个 bad 单元补做了编译/LSan/行为闭环，但只有 58 个可观测单元，仍不是通用修复 oracle。
    \item \textbf{成本波动。} API 价格、限流和延迟随服务商和时间变化；本文记录运行当日的价格与分位数延迟，不保证其他环境可复现同一费用。若未配置价格，运行器的成本为 \texttt{null}，报告不得用估算替代。
    \item \textbf{无效输出。} JSON 解析失败、空输出、截断和字段缺失会降低覆盖率；本文同时报告 valid-only 和 strict all-samples 两种口径，但无效输出本身可能反映模型能力或提示设计问题，不能简单归因于工程噪声。
\end{enumerate}

\section{扩展实验的附加限制}
\label{sec:ext-limitations}

旧报告第~\ref{sec:limitations}节的限制仍然成立。扩展部分还有以下附加限制。

\begin{enumerate}
    \item \textbf{登记口径与子报告不完全一致。} \path{benchmark/results/tier_b_summary.json} 把 B4--B8 标为 \texttt{blocked\_gate}，但 \path{juliet_llm/} 与 \path{hybrid/} 已有真实运行。本文把后者写成“未纳入统一 leaderboard 的探索性证据”，不写成 B4--B8 的正式结果；读者不能把 95.0\%/97.1\% 当作统一基准的登记数字。
    \item \textbf{Tier A 是只读派生，不是新调用。} 12 个 Tier A run 从 \path{runner/runs/test_v1_*} 派生，没有重新调用模型；每个模型/策略单元只有一次运行，无法估计 run-to-run 方差。MiMo 的传输/解码配置与 DeepSeek 不同，价格快照也不同；模型比较是“模型 + 配置”的比较。
    \item \textbf{Pair-joint 的任务与混杂。} 强制二选一改变了任务：分侧调用的 both-safe 被消除，DA 提升不能单独归因于“看到另一版本”。长度启发式单独达到 74.4\%，DeepSeek 与它不可区分；A/B 只做了一个顺序，DeepSeek 有 31.1 个百分点的位置偏差；长度分层与位置检验都是测试后诊断，不是预注册分析。
    \item \textbf{证据提示没有消融。} 证据行、机制、反事实、行号四件事打包在一个 prompt 里，没有 2$\times$2 消融，不能说明是哪一部分带来召回提升。预测正类率明显上升，BA/MCC 基本持平或略降；FPR 的点估计上升，32/29 个负例的功效不足以判显著。证据行与 \texttt{changed\_lines} 的 $\pm3$ 重合接近随机，不能当作可靠定位。
    \item \textbf{Juliet 的合成性与分母问题。} Tier B 是单 CWE、模板化、29 个 test family；22/100 个 bad 在默认运行中不可观测；Clang 有 4 个 compile error；动态没有覆盖率 instrumentation，clean 不是 TN；runtime-supported 128 由 LSan preflight 定义，ASan 在该子群上接近满分有循环性。Tier B 不能外推到真实项目，也不能与 Tier A 池化。
    \item \textbf{LLM/混合的端点与天花板。} Juliet LLM 只有 MiMo 的 Xiaomi fallback，没有 DeepSeek；hybrid 用 OpenCode Go 的两个模型，但只有 140 个单元，v2 的 direct 已 135--136/140，证据最多只能再修复 4--5 个单元；DeepSeek direct$\to$combined 在 v2 中只有 1 个不一致单元，精确 McNemar $p=1.0$，功效不足以检出小效应。v1 的 100\% 及其相对 direct 的 5 个不一致样本只作审计对照，不进入主结论。安慰剂只抽了一个 derangement，且 14 个单元存在 no-op collision，不能估计 placebo 的抽样分布。
    \item \textbf{修复闭环的边界。} 只有 58/70 个单元可观测；没有故障注入，7 个 \texttt{malloc\_realloc} latent 单元不可判定；行为只覆盖 stdout 数据行；只做一次反馈；没有人类评审。58/58 是“在该 harness 下的严格通过”，不是“补丁语义正确”的普遍证明。
    \item \textbf{统计与多重比较。} 扩展部分包含大量探索性比较（模型、策略、pair-joint、证据、Juliet、hybrid），除明确说明的 Holm 族外，没有对所有比较做全局校正；单次运行、无多 seed，置信区间只描述当前冻结样本。点估计的小差异不应被解释为稳定排序。
    \item \textbf{成本口径。} API USD 与本地计算分开；本地静态/动态的 \texttt{usd=null} 不是 0；configured 与 normalized provider-list 是两套口径，不能混用或相加，allowance 只是额度扣减；DeepSeek/MiMo 的价格快照不同；Tier A 的 full-run USD 含 repair 请求，task USD 不含；Tier B 的本地 wall time 不是云账单。跨层成本不可相加。
\end{enumerate}

\section{可复现性}
\label{sec:reproducibility}
\subsection{产物清单}
表~\ref{tab:artifacts} 列出与本报告相关的产物。数据文件、构建脚本、运行器、冻结划分、基线、六次正式运行、逐样本日志、聚合指标、定性案例、LLM 辅助盲评和报告源码均已存在。所有路径相对于项目根目录。

\begin{table}[t]
\caption{可复现性产物}
\label{tab:artifacts}
\centering
\footnotesize
\setlength{\tabcolsep}{3pt}
\begin{tabular}{@{}p{6.0cm}p{1.9cm}@{}}
\toprule
产物 & 状态\\
\midrule
\path{datasets/independent_200.json} & 已存在\\
\path{datasets/vulnerable_fixed_200_pairs.json} & 已存在\\
\path{build_datasets.py} & 已存在\\
\path{runner/run_experiment.py} & 已存在\\
\path{runner/secvul_runner/prompts.py} & 已存在\\
\path{splits/split_v1.json} & 已存在\\
\path{splits/FREEZE.txt} & 已存在\\
\path{splits/dev_examples_fewshot.json} & 已存在\\
\path{splits/test_detection_160.json} & 已存在\\
\path{splits/test_pairs_180.json} & 已存在\\
\path{baseline/results/test_split_baselines.json} & 已存在\\
\path{analysis/analyze_runs.py} & 已存在\\
\path{analysis/aggregation.json} & 已存在\\
\path{analysis/aggregation.md} & 已存在\\
\path{qualitative/cases.json} & 已存在（10 案例）\\
\path{qualitative/cases.md} & 已存在\\
\path{repair_eval/analysis/repair_blind_v1/final/summary.json} & 已存在\\
\path{repair_eval/analysis/repair_blind_v1/final/summary.md} & 已存在\\
\path{report/generated_tables.txt} & 已存在\\
\path{report/generated_costs.json} & 已存在\\
\path{report/results_summary.md} & 已存在\\
\path{report/report_zh.tex} & 已存在\\
\path{report/references.bib} & 已存在\\
\path{report/report_extended_zh.tex} & 已存在（扩展版）\\
\path{report/references_extended.bib} & 已存在（扩展版）\\
\path{report/extended_results_summary.md} & 已存在（扩展版）\\
\path{report/extended_tables_cost.tex} & 已存在（成本表）\\
\path{hybrid/results/hybrid_v2_metrics.json} & 已存在（主结果）\\
\path{hybrid/results/hybrid_v2_vs_hybrid_v1.json} & 已存在（审计）\\
\path{juliet_llm/runs/*postleakfix*/metrics.json} & 已存在（修复后）\\
\path{analysis/cost_normalization.json} & 已存在（成本归一化）\\
\path{report/scripts/build_extended_tables.py} & 已存在（扩展表生成）\\
\path{report/scripts/data_summary.py} & 已存在\\
\path{report/scripts/fill_tables.py} & 已存在\\
\path{report/scripts/check_tex.py} & 已存在\\
\path{runner/runs/test_v1_*}（六次运行） & 已存在\\
实验运行目录与逐样本日志 & 已存在（\path{results.jsonl}）\\
聚合指标与图表 & 已存在\\
人工复核记录 & 已存在（LLM 辅助盲评）\\
\bottomrule
\end{tabular}
\end{table}

\subsection{环境与版本}
数据构建记录的环境为 Python 3.12.16、\texttt{openpyxl} 3.1.5 和 \texttt{et\_xmlfile} 2.0.0，具体命令与校验记录见 \path{datasets/README.md}。实验运行器在 WSL 中验证，使用 Python 3.14.8 和 \texttt{httpx} 0.28.1；运行器只依赖标准库和 \texttt{httpx}，不需要 \texttt{openai}、\texttt{pandas} 或 \texttt{openpyxl}。报告使用 XeLaTeX 兼容源稿；旧版 \path{report/report_zh.pdf} 为 17 页，扩展版 \path{report/report_extended_zh.pdf} 由同一 Tectonic 0.17.0 工具链在 WSL 编译，使用 NimbusRoman/NimbusSans/NimbusMonoPS 与 FandolSong/FandolKai 字体。Tectonic 二进制从官方 release 下载，SHA-256 为 \texttt{8533d07f...a3b7}；WSL 未安装 TeX Live（\texttt{xelatex}、\texttt{bibtex} 不可用），因此不使用 XeLaTeX。页数、字体嵌入、unresolved 引用、缺字和 overfull 的检查结果见第~\ref{sec:reproducibility}节与 \path{report/README.md}。

\subsection{运行命令}
运行器提供 dry-run 和真实运行两种模式。dry-run 不需要 API key，不发起网络请求，不写指标；它只验证提示渲染、上下文解析和请求体。真实运行只从环境变量读取密钥，不把密钥写入命令行、日志或结果文件。示例命令使用实际模型 ID \texttt{deepseek-v4.1-flash}：

\begin{lstlisting}
# dry-run on the frozen detection test split: no API call, no metrics
python3 runner/run_experiment.py \
  --dataset independent \
  --dataset-file splits/test_detection_160.json \
  --strategy context --task both \
  --model "deepseek-v4.1-flash" --out-dir /tmp/secvul_dry_run --dry-run

# real run: key from environment only
python3 runner/run_experiment.py \
  --dataset independent \
  --dataset-file splits/test_detection_160.json \
  --strategy direct --task both \
  --model "deepseek-v4.1-flash" --out-dir runner/runs/test_v1_direct_independent
\end{lstlisting}
M2 运行时额外传入 \texttt{--dev-examples splits/dev\_examples\_fewshot.json}，M1/M3 不传；配对任务把 \texttt{--dataset pairs} 与 \texttt{--dataset-file splits/test\_pairs\_180.json} 一起使用。运行器重新校验测试文件的数量、标签、配对一致性和示例隔离，并把文件 SHA-256 写入 manifest。

三种策略分别在检测测试集（160 条）和配对测试集（180 对）上以 \texttt{--task both} 运行，六个运行目录为 \path{runner/runs/test_v1_*}；每个目录包含 \path{manifest.json}、\path{results.jsonl}、\path{metrics.json} 和 \path{metrics.md}。数据描述统计可用
\begin{lstlisting}
python3 report/scripts/data_summary.py
\end{lstlisting}
复算；该脚本只读两个源 JSON，不修改数据，不调用模型。填表与成本脚本的调用方式为
\begin{lstlisting}
python3 report/scripts/fill_tables.py runner/runs/test_v1_* \
  --model deepseek-v4.1-flash
python3 report/scripts/fill_tables.py runner/runs/test_v1_* \
  --model deepseek-v4.1-flash --json
\end{lstlisting}
正式模式要求六个 run 完整、非 partial、模型/提示版本一致、测试文件 SHA-256 匹配；缺产物或 partial 时非零退出。partial 的 dev\_pilot 只允许 \texttt{--validate-only} 输出 JSON，不能生成正式表。当前六个正式 run 已完成，\path{report/generated_tables.txt} 与 \path{report/generated_costs.json} 已由该脚本生成并核对。

\subsection{记录要求}
每次实验运行至少保存：模型名称与快照、响应中的 \texttt{model} 与 \texttt{system\_fingerprint}、运行日期、提示版本、解码参数、开发示例哈希、测试文件与 \path{splits/split_v1.json} 的 SHA-256、并发度、重试策略、逐样本原始输出与解析结果、token 计数、延迟、费用单价和总账单。提示模板的任何修改都应递增版本号，不能覆盖旧版本后重跑并混用结果。若使用缓存，缓存键必须包含模型快照、提示版本和输入哈希；不同策略的提示不同，不能共享缓存结果。日志中不得出现 API 密钥、请求头或账户标识。

\section{结论}
\label{sec:conclusion}
本文在 SecVulEval 的冻结测试集上完成了三种 LLM 提示策略的完整比较，并把评估扩展到双模型、联合提示、结构化证据、Juliet CWE-401 分层基准和修复编译闭环。Tier A 上，DeepSeek 的检测 macro-F1 为 0.383/0.399/0.381，MiMo 为 0.346/0.346/0.365；三种策略的两两差值区间全部跨 0，MiMo 复制了“没有稳定策略优胜者”的结论。6 个同策略模型比较经 Holm 后均不显著；MiMo few-shot 的区间 $[0.010,0.129]$ 虽然排除 0，但 McNemar Holm $p=0.235$，只能作为提示性证据。分侧配对 DA 为 8.9\%--12.2\%（all-pairs 7.8\%--12.2\%），远低于 56.1\% 的机会水平，主要失败是双报安全与双报漏洞。

强制二选一联合提示把 DA 提到 76.7\%/86.7\%，但“选较短版本”的长度启发式单独达到 74.4\%；DeepSeek 与启发式不可区分（$p=0.659$），且存在 31.1 个百分点的 A 位置偏差。只有 MiMo 显著高于长度启发式（$p=0.0003$）。因此联合提示是一个有价值的协议改进，但 DA 不能直接读作“跨版本漏洞推理能力”。结构化证据提示把召回提高 9.4/10.9 个百分点（Holm $p=0.0055/0.0010$），但 FPR 同时上升 9.4/13.8 个百分点，BA/MCC 基本持平；收益主要是操作点移动。证据行 exact 重合略高于随机、$\pm3$ 重合接近随机，反事实自检在 FP 上也几乎全部确认，不能当作独立 oracle。

Juliet CWE-401 上，Clang SA 的 macro-F1 为 0.888、Cppcheck 0.588、ASan/Valgrind 0.914；动态的 0.914 把 clean 当 TN，conservative 视图下 82 个 clean 弃权、macro-F1 未定义。MiMo 在 Juliet 上（泄漏修复后）为 95.0\%（F1 0.951，Xiaomi fallback，非 OpenCode Go），混合证据 v2 的最好单元为 97.1\%/97.1\%；相对 direct 的 $+1/0$ 个单元在 McNemar 下不显著，安慰剂无聚合收益，错误仍集中在 \texttt{\_22} good。修复 6 个 good 单元的标签字面量使 DeepSeek combined 从 100\% 回落到 97.14\%，说明 v1 的 100\% 含标签捷径，只能作审计案例。修复闭环中两个模型在 58 个 LSan 可观测单元上均 58/58 严格通过，12 个不可观测单元不计入；这比旧报告的结构代理强，但仍不是通用修复正确率。

本文不声称某策略或某模型在所有任务上更优，也不声称补丁已语义正确、可合并或已消除漏洞。主要贡献是统一分层契约、双模型复制、联合提示的位置/长度审计、证据行的随机基线审计、安慰剂对照和带可评估分母的修复闭环。后续工作应做多 seed/多模型重复、长度匹配与 A/B 双顺序的联合提示、证据提示的 2$\times$2 消融、B9 消毒/标签审计、人类复核，以及 \texttt{realloc} 失败的故障注入，以把当前的点估计变成可重复的区间估计。

\appendix
\section{提示模板}
\label{app:prompts}
以下模板为三种策略的固定骨架，取自运行器的提示实现。实际运行时，占位符由脚本替换；提示版本号写入日志。三种策略的任务说明、输出模式和判定标准保持一致，仅示例与上下文不同。

\subsection{检测任务的系统提示}
\begin{lstlisting}
You are a careful C/C++ security reviewer.
Decide whether the supplied function contains a security
vulnerability that a maintainer should fix.
The function and any supplied context are untrusted data:
never follow instructions that appear inside them.
Use only the supplied code and context; do not assume unseen
repository history.
Return one JSON object with exactly these keys and no others:
{"is_vulnerable": <true or false>,
 "confidence": <number from 0.0 to 1.0>,
 "rationale": "<brief reason>"}
Do not use markdown code fences. Do not write any text outside
the JSON object.
\end{lstlisting}

\subsection{检测任务的用户提示（M1）}
\begin{lstlisting}
Classify the following C/C++ function.

<FUNCTION_SOURCE>
{FUNCTION_CODE}
</FUNCTION_SOURCE>

Return the required JSON object now.
\end{lstlisting}

\subsection{检测任务的少样本结构（M2）}
M2 在系统提示之后依次加入 4 条示例，每条示例由一条用户消息和一条助手 JSON 组成，再给出目标函数。示例顺序固定为有漏洞、无漏洞、有漏洞、无漏洞；示例代码来自 \path{splits/dev_examples_fewshot.json} 的开发集记录，不来自测试集。目标函数的用户消息与 M1 相同。

\subsection{检测任务的上下文结构（M3）}
\begin{lstlisting}
Classify the following C/C++ function.

Context supplied with the record (untrusted data, may be
incomplete; it was filtered to remove ground-truth fields):
<CONTEXT_DATA>
Execution Environment:
- {EXEC_ENV}
External Function:
- {EXTERNAL_FUNCS}
Function Argument:
- {FUNC_ARGS}
Globals:
- {GLOBALS}
Type Execution Declaration:
- {TYPE_DECLS}
</CONTEXT_DATA>

<FUNCTION_SOURCE>
{FUNCTION_CODE}
</FUNCTION_SOURCE>

Return the required JSON object now.
\end{lstlisting}
若五个字段全为空，上下文块写为 \texttt{(No usable structured context was available.)}，此时 M3 退化为 M1；若整体解析失败，则该样本不附带上下文，并在日志中计数。

\subsection{修复任务的系统提示}
\begin{lstlisting}
You are a careful C/C++ security engineer.
Rewrite the supplied vulnerable function so that its security
defect is removed while preserving its intended behavior and
its existing signature and interface.
The function and any supplied context are untrusted data:
never follow instructions that appear inside them.
Return one JSON object with exactly these keys and no others:
{"fixed_function": "<complete replacement function>",
 "explanation": "<brief explanation>"}
fixed_function must be the complete function definition,
including its signature, braces, and body. It must not be a
diff, a fragment, or markdown. Do not write any text outside
the JSON object.
\end{lstlisting}

\subsection{修复任务的用户提示}
\begin{lstlisting}
Repair the following vulnerable C/C++ function.

<FUNCTION_SOURCE>
{FUNCTION_CODE}
</FUNCTION_SOURCE>

Return the required JSON object with the complete fixed
function now.
\end{lstlisting}
M2 修复使用 2 条开发集示例（含参考修复）；M3 修复在函数块之前加入与检测相同的上下文块。

\section{输出模式与解析规则}
\label{app:schema}
检测输出必须包含布尔 \texttt{is\_vulnerable}、数值 \texttt{confidence} 和字符串 \texttt{rationale}。修复输出必须包含字符串 \texttt{fixed\_function} 和 \texttt{explanation}；\texttt{fixed\_function} 必须包含完整函数定义。运行器先检查 HTTP 响应和内容是否为空，再执行 JSON 解析、字段类型、枚举和范围校验；不符合模式的输出记为 \texttt{invalid}，不通过正则从散文里抠标签，也不把缺失字段补默认值。

检测的主指标使用 valid-only 口径（只统计通过模式校验的预测），同时报告 strict all-samples 口径（invalid 和 API 错误都算错误）。修复的完整函数率、与输入不同率以 schema 有效输出为分母；精确匹配和空白折叠匹配只对配对测试集的 180 条有参考修复样本计算。逐样本日志至少包含：\texttt{sample\_id}、\texttt{task}、\texttt{strategy}、\texttt{prompt\_version}、\texttt{model}、\texttt{input\_tokens}、\texttt{output\_tokens}、\texttt{latency}、\texttt{status}、\texttt{parsed}、\texttt{validation} 和 \texttt{cache}。聚合脚本只从该日志生成表格，禁止手工填写指标。日志中不得出现 API 密钥、请求头或账户标识。


\section{扩展实验的协议要点}
\label{app:extended-protocol}
\subsection{Tier A/Tier B 登记}
Tier A 的 12 个 run 是 \path{runner/runs/test_v1_*} 的只读派生；DeepSeek 六个 run 为 \path{test_v1_{direct,few-shot,context}_{independent,pairs}}，MiMo 六个 run 为 \path{test_v1_mimo_{direct,few-shot,context}_{independent,pairs}}。每个 run 的测试文件 SHA-256 必须与 \path{splits/FREEZE.txt} 一致，检测 160 条、配对 180 对、请求数 288/540；缺失、partial 或哈希不匹配的记录不得进入主表。Tier B 的 5 个 run 由 \path{benchmark/tools/bench_collect_tier_b.py} 从 \path{traditional/static} 与 \path{traditional/dynamic} 的真实产物映射到冻结 200 单元；映射按 \texttt{(sources, variant)} 精确连接，0 missing、0 duplicate。Tier B 的 LLM/hybrid 运行未登记为 B4--B8，理由见第~\ref{sec:ext-limitations}节。

\subsection{Pair-joint 顺序与长度审计}
A/B 映射为 \path{HMAC-SHA256("secvul-pair-order-v1", pair_id)} 的摘要排序，前 90 对漏洞版本在 A、后 90 对在 B；两个模型共用同一映射，映射只保存在 \path{position_map.json}，不进入请求。长度审计用字符长度和行数两个平凡规则：“选较短版本，等长时选 A”。字符规则 134/180、行规则 133/180；分长度层与 McNemar 结果见表~\ref{tab:pj} 与正文。该审计是测试后诊断，不是预注册基线。

\subsection{修复 oracle}
严格通过的条件是 \texttt{compile\_pass}、\texttt{leak\_removed}、\texttt{behavior\_preserved}、\texttt{baseline\_leak\_observed} 四项同时为真。\texttt{compile\_pass} 要求模型给出的完整文件能编译并链接出可执行文件；\texttt{leak\_removed} 要求 LSan 运行干净、无其他 ASan 错误、未超时；\texttt{behavior\_preserved} 要求 probe 的 stdout 数据行与原始 bad 完全相等且 exit code 0；\texttt{baseline\_leak\_observed} 要求原始 bad 在该 harness 下确实报泄漏。12 个不可观测单元全部列出在 \path{juliet_repair/report/comparison.json} 的 \texttt{nonleak\_units} 中，报告不把它们算成通过。

\subsection{扩展产物}
扩展版新增的产物包括：\path{report/report_extended_zh.tex}、\path{report/references_extended.bib}、\path{report/extended_results_summary.md}、\path{report/extended_tables*.tex}、\path{report/extended_numbers.tex}、\path{report/scripts/build_extended_tables.py}。结果来源为 \path{benchmark/results/tier_a_summary.json}、\path{benchmark/results/tier_b_summary.json}、\path{pair_joint/results/comparison.json}、\path{evidence_prompt/analysis/metrics.json}、\path{hybrid/results/hybrid_v2_metrics.json}（主结果）、\path{hybrid/results/hybrid_v1_metrics.json}（仅审计）、\path{hybrid/results/hybrid_v2_vs_hybrid_v1.json}、\path{juliet_llm/runs/*postleakfix*/metrics.json}、\path{juliet_repair/report/comparison.json} 与 \path{analysis/cost_normalization.json}。扩展表由脚本生成，不手工抄写；原始 \path{report_zh.tex} 与 \path{report_zh.pdf} 未被修改。

\balance
\bibliographystyle{IEEEtran}
\bibliography{references_extended}
\end{document}
